\documentclass[11pt]{article}%
\usepackage{amsmath}
%\usepackage{appendix}
\usepackage{amsfonts}
\usepackage{amssymb}
\usepackage{booktabs} 
\usepackage{pdflscape}
\usepackage{graphicx}
\usepackage{caption}
\usepackage{subfig}
\usepackage{geometry}
\usepackage[sidewaysfigure]{rotating}
\usepackage{natbib}
\bibpunct{(}{)}{;}{a}{,}{,}
\usepackage[doublespacing]{setspace}%
\setcounter{MaxMatrixCols}{30}

\usepackage{amsthm}
\newtheorem{assumption}{Assumption}
\newtheorem{definition}{Definition}
\newtheorem{proposition}{Proposition}
\newtheorem{corollary}{Corollary}
\newtheorem{lemma}{Lemma}
\newtheorem{theorem}{Theorem}
\renewcommand{\qedsymbol}{\textit{Q.E.D.}}

\providecommand{\U}[1]{\protect\rule{.1in}{.1in}}

\geometry{left=2.1cm,right=2.1cm,top=2.1cm,bottom=2.1cm}

\usepackage{hyperref}
\hypersetup{
  pdftitle={OLG Models},    % title
  pdfauthor={Robert Kirkby},     % author
  %hyperfootnotes=false,
  %hyperindex=true,
  %hyperfigures=true,
  colorlinks=true,                % makes links a coloured word, rather than putting coloured boxes around them
  linkcolor=blue,                % makes internal links black
  citecolor=blue,                % makes citations black
  urlcolor=blue,                 % color of external links
  breaklinks=true                 % allow links to be broken across multiple lines
}

\newcommand{\rob}[1]{{\color{red}\bf #1}}


\begin{document}

\title{An Introduction to OLG Models}
\author{Robert Kirkby \\ Victoria University of Wellington
\\ \href{https://doi.org/10.5281/zenodo.8136815}{https://doi.org/10.5281/zenodo.8136815}
% \thanks{The authors would like to thank ...}
}
\date{\today \\ }
\maketitle

\vspace{-10mm}
\begin{abstract}
Explores and explains some basic OLG models. Starts with a deterministic OLG, showing how to define and compute the stationary general equilibrium. After this basic model it is shown how to add age-dependent parameters, survival probabilities (the possibility of dying before the final period), assets, a pension system, progressive taxation and government spending. We then move to stochastic OLGs by introducing idiosyncratic shocks, both persistent and transitory. Later models include: an OLG model in which there are married couple households, single male households, single female households; an OLG model with both heterogeneous households and heterogenous firms; an OLG model with a search-and-matching labor market. This is done over the course of a number of models to illustrate and explain clearly one-by-one how these parts are added to the model and implemented in the codes for each model.

It is recommended to first read the Intro to Life-Cycle Models: \href{https://www.vfitoolkit.com/updates-blog/2021/an-introduction-to-life-cycle-models/}{vfitoolkit.com/updates-blog/2021/an-introduction-to-life-cycle-models/}
\end{abstract}

Keywords: Over-lapping generations; Fiscal; General equilibrium. \\
\small

\thispagestyle{empty}

\newpage

\tableofcontents

\newpage

\newpage

\section{Introduction}
Explores and explains some overlapping-generations (OLG) models. Starts with a deterministic OLG, showing how to define and compute the stationary equilibrium. After this basic model it is shown how to add age-dependent parameters, survival probabilities (the possibility of dying before the final period), assets, a pension system, progressive taxation and government spending. Idiosyncratic shocks are then introduced leading us to stochastic OLGs. By the end it is solving an OLG model with heterogeneous households and heterogeneous firms, in which there are households facing genuinely different problems for married couples, single males, and single females. This pdf explains the models one-by-one, and the codes implementing each of them are all on github. We make use of the VFI Toolkit (vfitoolkit.com) and you will need Matlab with a compatible GPU.

If you are not already familiar with finite-horizon life-cycle models and how to solve them, I recommended you learn that first although you could skip it if you want, see: \href{https://www.vfitoolkit.com/updates-blog/2021/an-introduction-to-life-cycle-models/}{\textcolor{blue}{VFI Toolkit: An Introduction to Life-Cycle Models}}

If you are wanting to solve any kind of OLG model with one endogenous state, you probably will find everything you need by the end of this introduction to OLG models. If you want to solve advanced models with multiple endogenous states this intro will help you understand all the steps and how to set things up, but you will need to go and learn more sophisticated (but less easily generalized) algorithms needed to solve such models.

The grids used for this introduction are on the smaller side so as to make the models more easily solved and accessible in terms of computational requirements. As a result you likely want to use larger grids in practice. Most models require a GPU with 8gb of gpu memory (GDDR memory) that works with Matlab. The first five models (the deterministic OLG) will likely run with much less (I guess 2gb?). The last few models, especially those with married couples will need 40gb of gpu memory.

By the end of this introduction you will hopefully have a good idea of how to put together OLG models. If you are wanting to solve one-asset models then you can probably just use the toolkit as is. If you want to solve models with more assets you will likely need to go and learn some more advanced algorithms for solving OLG models (in particular the value function problem).

A follow-up 'Introduction to OLG Transition Paths' should hopefully be forthcoming in a year or so.

\newpage
\section{The Models}
We first build a deterministic OLG model, adding assets and a government that uses progressive taxes and runs a balanced budget, over the course of the first five models. Then in OLG Model 6 we make switch to a stochastic OLG model (idiosyncratic, but no aggregate, shocks), the change is very simple to make. OLG Model 7 then simply takes this model and shows how to calculate a variety of things related to it such as aggregates, life-cycle profiles, simulating panel data, and welfare.

For OLG models 1-7 we will simply solve them in per-capita terms without really discussing this issue (simply as this is the mathematically and conceptually easiest way to approach the models). In OLG model 8, when we introduce deterministic economic/productivity growth we will also discuss how to deal with population growth in terms of solving for aggregate quantities, rather than just per-capita quantities as will have been implicit up until then. In practice, most OLG models in the literature are solved in per-capita terms.

In OLG model 10 we introduce 'permament types' to solve an OLG Model with a fixed-effect in labor productivity. We then use this to build up a model in which we have married couples, single males, and single females; differences being not just things like different parameters, but different exogenous shocks, and even different numbers of decision variables (the married couple has two labor supply decisions, one for each spouse).

In OLG model 14 we introduce heterogeneous firms, to have a model with both OLG heterogeneous households and heterogeneous firms.

In OLG model 16 we introduce a search-and-matching labor market.

In the codes we always write the general equilibrium conditions as an expression that will evaluate to zero when general equilibrium holds.

The basic approach to solving these OLG models is very easy to use. Appendix A gives a brief description of the concept of 'solving' to find the general equilibrium involves. Appendix B shows how to speed it up by changing the algorithm for finding the general equilibrium, which requires giving just a little bit more information, it also describes how to instead use a slower but more robust algorithm to find the general equilibrium.


\newpage
\section{Building up OLG Models}


\subsection{OLG Model 1: Consumption-Leisure in general equilibrium OLG}
Households live for J periods and solve the life-cycle problem,
\begin{eqnarray*}
 \sum_{j=1}^J \beta^{j-1} && \left[ \frac{c_j^{1-\sigma}}{1-\sigma} - \psi \frac{h_j^{1+\eta}}{1+\eta} \right] \\
&&\text{if working, }j<=Jr: c_j=(1-\tau) w h_j \\
&&\text{if retired, }j\geq Jr: c_j=pension \\
&& 0\leq h_j \leq 1
\end{eqnarray*}
where $j$ indexes age, $c_j$ is consumption, $h_j$ is hours worked,\footnote{$h$ is actually 'fraction of time worked', but I will call it hours worked throughout simply as it is easier to think of in this way. $0\leq h \leq 1$ is showing us that this 'fraction of time worked' is from 0 to 1.} $w$ is the wage per hour worked. So a household that works $h_j$ hours receives an hourly wage of $w$, giving them a pre-tax labor income of $w h_j$. Notice that all the parameters are just a single number. $\tau$ is a tax rate applied to labor income (here set to zero, in OLG Model 2 it will be used to fund the pensions). $pension$ is a pension received by retirees (notice that we do not allow retirees a choice about whether or not to work). $Jr$ is the retirement age.

We can rewrite this as a value function problem,
\begin{eqnarray*}
 V(j)&=& \max_{c,h} \frac{c^{1-\sigma}}{1-\sigma} - \psi \frac{h^{1+\eta}}{1+\eta}  + \beta V(j+1) \\
&&\text{if working, }j<=Jr: c=(1-\tau) w h \\
&&\text{if retired, }j\geq Jr: c=pension \\
&& 0\leq h \leq 1
\end{eqnarray*}
notice that the value function $V$ depends on the age $j$, so households make different decisions (different optimal policies) at different ages. We assume $V(J+1)=0$.\footnote{$V(J+1)$ is effectively the utility received upon dying.} Notice also from the budget constraint that once the household chooses one of $c$ or $h$, it is implicitly choosing the other. The codes take advantage of this to just have one decision of $h$.

The other part we need to add to this model is the firm side of the economy. We will just use the simplest implementation possible, which is an aggregate production function that has decreasing returns to labor,\footnote{A representative firm is possible with constant returns to scale production functions as the distribution across individual firms of production inputs is irrelevant. This is not true when, as here, there are decreasing returns to scale; the distribution of production inputs across firms matters and a representative firm is not possible. Since we plan to switch to a Cobb-Douglas production function with constant returns to scale in the following models we will simply ignore this issue here. For now we simply use this production function and ignore that it cannot represent a market aggregation of individual firms.}
\begin{equation}
 Y=L^{1-\alpha}
\end{equation}

We assume there are perfectly competitive labor markets, and so in general equilibrium it follows that the wage is equal to the marginal product of labor, which gives us the general equilibrium condition,
\begin{equation}
 w=(1-\alpha)L^{-\alpha}
\end{equation}
The aggregate labor $L$, is calculated by integrating over the hours worked by each of the individual households. We use $\mu$ to denote the distribution of agents, and so $L= \int h d\mu$; we are summing up the hours worked of each household, and $\mu$ is providing the 'weight' for each household type. In this first model $\mu$ will just be the fractions of agents of each age, but it will become more meaningful in later models.

Any OLG model requires us to define how many people there are of each age, which the codes call $mewj$, which is the marginal distribution of $\mu$ of households/agents over age $j$, and to define the state of people when they are born, which the codes call $jequaloneDist$. We will discuss these in later models as they are trivial and uninteresting in this first model where we just set $mewj$ so that there are an equal fraction of the population in each age.

The definition of a stationary equilibrium for this model is,
\begin{definition}
 A stationary equilibrium is price $w$,
 \begin{enumerate}
  \item Given prices ($w$), the household value function, $V$, and related policy function, $policy$, solve the household problem.
  \item The agents distribution evolves according to,
   \\  $\mu(j+1)=\int \mu(j), \; j=1,2,...,J$, and $\mu(1)=jequaloneDist$
  \item The aggregates are based on household policies and agent distribution,
   \\  $L=\int policy_h d\mu$
  \item Markets clear: $w=(1-\alpha)L^{-\alpha}$
 \end{enumerate}
\end{definition}
Notice this first model has no endogenous state in the household problem and so the agent distribution is actually independent of the household policies.

Note also that is is common to include $V$, $policy$, and $L$ in the definition of the stationary equilibrium. I do not do so here mainly to save notation, and because these can anyway always be calculated given the prices and government policies, which is how the codes will work (we find the equilibrium prices and government policies, and then from these can calculate things like $V$, $policy$ and $L$ as desired).

When we come to write the codes we will rewrite the market clearance condition, which is the only general equilibrium condition, as an expression that equals zero in equilibrium. So we write it as $w-(1-\alpha)L^{-\alpha}$ in the codes.

\subsection{OLG Model 2: Use tax to fund pensions}
The only difference from OLG Model 1 is that we now use the tax, $\tau$ (which was previously equal to zero) to balance the pensions expenditures. This involves adding a general equilibrium condition to determine the tax rate such that the tax revenue raised is equal to the total expenditure on pensions. This shows how we add general equilibrium conditions, and also illustrates how government budget conditions are treated just like any other general equilibrium condition.

Since the household problem is unchanged ($\tau$ was already in OLG Model 1, just that it was set equal to zero), we skip rewriting it here. Same goes for the firm side of the economy and the labor market clearance.

The only other thing we have is the pension system, we will require that the pension system runs a balanced budget, so total revenues raised by the tax rate $\tau$ must equal total spending on the pension. We will require that the pension system has a balanced budget: revenues from the payroll tax equal total spending on pensions for the retirees. In the current model this is easy, we just require $pension*((J-Jr)/J)=\tau*w*L$, where the left-hand side is the total spending on pensions and the right-hand side is the total revenue from payroll tax $\tau$.\footnote{Key to this equation is that we know that $(J-Jr)/J$ is the fraction of the population that is retired and receiving pensions, and that the payroll tax is a flat-rate so we can just multiply it by the tax base $w*L$.}

The definition of a stationary equilibrium for this model is,
\begin{definition}
 A stationary equilibrium is price $w$, and government policies $\tau$ and $pension$ such that
 \begin{enumerate}
  \item Given prices ($w$) and government policies ($\tau$, $pension$), the household value function, $V$, and related policy function, $policy$, solve the household problem.
  \item The agents distribution evolves according to,
   \\  $\mu(j+1)=\int \mu(j), \; j=1,2,...,J$, and $\mu(1)=jequaloneDist$
  \item The aggregates are based on household policies and agent distribution,
   \\  $L=\int policy_h d\mu$
  \item Markets clear: $w=(1-\alpha)L^{-\alpha}$
  \item Pension system balance: $pension*((J-Jr)/J)=\tau*w*L$
 \end{enumerate}
\end{definition}
Notice this model has no endogenous state in the household problem and so the agent distribution is actually independent of the household policies.

Note also that is is common to include $V$, $policy$, and $L$ in the definition of the stationary equilibrium. I do not do so here mainly to save notation, and because these can anyway always be calculated given the prices and government policies, which is how the codes will work (we find the equilibrium prices and government policies, and then from these can calculate things like $V$, $policy$ and $L$ as desired).

\subsection{OLG Model 3: Dependence on Age}
We make three extensions to the second model. Our first change is to allow wages to vary with age: we include a deterministic life cycle profile of labor efficiency unit, $\kappa_j$, intended to capture the empirical fact that average wages increase with age until something like 45-55 years old, and then decrease until retirement. Our second change is to include a 'survival probability', so that as people get older they become less likely to survive until the next period/age. Note that this will mean there are less elderly people giving us a more realistic age-demographic. The third is to include population growth at a rate $n$, which also helps us capture population age-demographics.

We will make a small change to the pension system budget balance, rather than fixing pensions and choosing $\tau$ to make the budget balance as in OLG Model 2, we will fix $\tau$ and choose pensions to make the budget balance. Note that the general equilibrium condition relating to the pension system budget balance is unchanged, we just change the name for which variables are to be determined in general equilibrium. We do this solely to illustrate how easily it can be done.

The households problem now reflects these changes,
\begin{eqnarray*}
 V(j)&=& \frac{c^{1-\sigma}}{1-\sigma} - \psi \frac{h^{1+\eta}}{1+\eta}  + \beta s_j V(j+1) \\
&&\text{if working, }j<=Jr: c=(1-\tau) w \kappa_j h \\
&&\text{if retired, }j\geq Jr: c=pension \\
&& 0\leq h \leq 1
\end{eqnarray*}
The 'age-conditional survival probability' $s_j$ ---the probability of surviving to age $j+1$ given agent is alive at age $j$--- enters alongside the 'discount factor', $\beta$, reflecting that agents only care about next period if they are alive next period. The deterministic labor efficiency age-profile, $\kappa_j$, enters into the earnings $w \kappa_j h_j$, which also changes the interpretation of $w$ which is now the 'wage per efficiency unit hour', rather than just 'wage per hour'; a household's wage per hour is $w \kappa_j$.

There are two important changes to other aspects of the model. The first is that the total labor supply is no longer just adding up hours worked across the households, it now includes the efficiency units, and so the total labor supply is now measured in efficiency units. This gives us,
\begin{equation*}
 L=\int \kappa_j h d\mu
\end{equation*}
We can also define the total 'hours worked' as $H=\int h d\mu$, note that in this model $h$ is 'fraction of time worked' rather than 'hours', I call it 'hours worked' for convenience.

The second change is that because people have a chance of dying we need to alter how we keep track of how many people there are of each age. We do this using in the code using $mewj$, which is the marginal distribution of $\mu$ over age. If we started with a population mass of 1 at age one, so $\mu^1=1$, then the survival probabilities mean that  $\mu^{j+1}=s_j \mu^j$ for $j=1,2,...,J$. We also have population growth at rate $n$ ($n$ is a percentage as a fraction so, e.g., $n=0.02$ represents 2\% annual population growth rate) so if the population that was age one yesterday equals $1$, then the population that is age one today equals $1*(1+n)$ ---it has grown at rate $n$--- and the population that is age two today is the same mass of $1$ that was age one yesterday (ignoring survival probabilities for the moment); rearranging this we get that if age one today is mass $1$, then age two today is mass $1/(1+n)$. If we apply this logic to the formula we had we still have $\mu^1=1$, but now $\mu^{j+1}=s_j \mu^j/(1+n)$ for $j=1,2,...,J$, and this is the formula used in the codes. We want to keep the total population equal to one to make the model easier to use, and so we then normalize the total population to one.

Some model output that is of interest with this second model is the 'life-cycle profile' of mean hours worked. These are graphs with age on the horiontal axis that show average hours worked for each age. Analagously the life-cycle profile for mean effective labor supply shows average effective labor supply, $\kappa_j h$, for each age. A second model output of interest in this second model is the 'age-demographic', which we here draw as a demographic pyramid.

Notice that while the requirement that the pension system has a balanced budget is conceptually unchanged we do have to modify the exact equation because $(J-Jr)/J$ is no longer the fraction of the population that is retired because of the probability of death/survival. We could just calculate the fraction (it can be directly calculated from the $mewj$) prior to solving the model and put this into our equation. But we will instead take advantage of just how easy it is to implement using the codes and add it as an aggregate to be calculated. This will lead to a minor change to the pension system balance equation.

The main change to the definition of stationary equilibrium is the equation covering how the agents distribution evolves. Note that the value function problem is different, and so $V$ and $policy$ are different, but we don't need to change how we have written this in the equilibrium definition.

The definition of a stationary equilibrium for this model is,
\begin{definition}
 A stationary equilibrium is price $w$, and government policies $\tau$ and $pension$ such that
 \begin{enumerate}
  \item Given prices ($w$) and government policies ($\tau$, $pension$), the household value function, $V$, and related policy function, $policy$, solve the household problem.
  \item The agents distribution evolves according to,
   \\  $\mu(j+1)=\int \frac{\mu^{j+1}}{\mu^j} \mu(j), \; j=1,2,...,J,\text{ and }\mu(1)=jequaloneDist$
  \item The aggregates are based on household policies and agent distribution,
   \\ $L=\int \kappa_j policy_h d\mu$
   \\ $PensionSpending=\int pension \mathbb{I}_{j\geq Jr} d\mu$
  \item Markets clear: $w=(1-\alpha)L^{-\alpha}$
   \\ Pension system balance: $PensionSpending=\tau*w*L$
 \end{enumerate}
\end{definition}
\noindent Note the difference between $\mu(j)$ which is the population distribution at age $j$, and $\mu^j$ which is the marginal population distribution at age $j$; in the current model they are identical, but from the next model on they are different. $\mathbb{I}_{j\geq Jr}$ is the indicator function for $j\geq Jr$; that is a function that takes the value 1 when $j\geq Jr$, and zero otherwise.\footnote{More generally, the indicator function $\mathbb{I}_S$ for some condition $S$, is a function that takes a value of 1 when $S$ is true, and $0$ otherwise. It can also/alternatively be defined with $S$ being a set, in which case it takes a value of 1 when in the set $S$, and $0$ outside the set $S$.} Notice then that $\mathbb{I}_{j\geq Jr}$ is taking a value of 0 for working age, and 1 for retirement.


\subsection{OLG Model 4: Assets and Aggregate Physical Capital}
We make the big change of adding assets to the household problem, and these household assets add up to aggregate physical capital which appears in the production function of the representative firm. Because we use a Cobb-Douglas production function the wage and interest rate are isomorphic, that is, one is just a simple function of the other, and vice-versa. We could therefore write the model in terms of either $r$ or $w$, and we switch to writing the model in terms of $r$ simply to follow what is standard in the literature. Because people can die accidentally (due to the survival probabilities already introduced) and have assets when they do we need to do something so these assets don't just 'disappear'; we will redistribute all these assets left by those who die, called 'accidental bequests', as a lump-sum payment to surviving agents, denoted $AccidentBeq$. Empirically, people who die often leave behind substantial assets, whereas people in our model will aim to die with zero assets (the survival probabilities mean they won't be exactly zero, but they will be small), and so we will introduce a 'warm glow of bequest' to make households keep assets until they die; essentially we just give them utility for having assets when they die.

Begin with the household problem which now includes assets, $a$, and which we first write as a sequence problem with a J period life-cycle problem
\begin{eqnarray*}
 \sum_{j=1}^J \beta^{j-1} \Pi_{i=1}^{j-1}s_i && \left[ \frac{c_j^{1-\sigma}}{1-\sigma} - \psi \frac{h_j^{1+\eta}}{1+\eta} \right] \\
&&\text{if working, }j<=Jr: c_j+a_j=(1-\tau) w h_j + (1+r) a_{j-1} +(1+r)AccidentBeq \\
&&\text{if retired, }j\geq Jr: c_j+a_j=pension+ (1+r) a_{j-1} +(1+r)AccidentBeq + \mathbb{I}_{j=J} warmglow(a') \\
&& 0\leq h_j \leq 1
\end{eqnarray*}
The important thing when we change to value function notation is that the household assets, $a$, is now an 'endogenous state' of the value function (so decisions/optimal policy depend on this state),
\begin{eqnarray*}
 V(a,j)&=& \max_{c,h,a'} \frac{c^{1-\sigma}}{1-\sigma} - \psi \frac{h^{1+\eta}}{1+\eta}  + \beta s_j V(a',j+1) \\
&&\text{if working, }j<=Jr: c+a'=(1-\tau) w \kappa_j h + (1+r) a  +(1+r)AccidentBeq\\
&&\text{if retired, }j\geq Jr: c+a'=pension+ (1+r) a  +(1+r)AccidentBeq +\mathbb{I}_{j=J} warmglow(a') \\
&& 0\leq h \leq 1
\end{eqnarray*}
So now the value function has one endogenous state, $a$, and the age state $j$, and the household needs policies for both $a'$ and $h$.

The warm glow of bequests, is $\mathbb{I}_{j=J} warmglow(a')$, the first term is and indicator for final periods and is telling us that we only get the warm glow when we die (in last period) and the second part is telling us that the warm glow is a function of the assets that we choose, $a'$, in the final period and so are the assets we have at the end of the final period when we die. We will give the warm glow of bequests the functional form $warmglow(a')=wg_1 \frac{(1+a'/wg_2)^{1-wg_3}}{1-wg_3}$, where $wg_1$ controls the (relative) importance of bequests, $wg_2 \geq 1$ controls the degree to which bequests are a luxury good ($wg_2=1$ would be a normal good, larger $wg_2$ makes bequests more luxury-like so that wealthy households leave proportionally larger bequests), and $wg_3$ is the curvature parameter.\footnote{Notice that you only get the warm-glow if you die at the end of the final period, not if you die in any of the previous periods due to the survival probabilities. We could instead use $\beta (1-s_j)warmglow(a')$ to make it so we get the warm glow of bequest if household dies at any age; note that in this case we must set $s_J=0$. Philosophically this alternative seems nicer, but this is irrelevant in a quantitative model. I am not aware of anything establishing which of the two performs better empirically. I do not use the alternative here solely because it is slightly more subtle to explain.}\textsuperscript{,}\footnote{Notice that the warm-glow function has the same CRRA-style curvature as the utility of consumption (in the codes we set $wg_3=\sigma$). This gives it the same concave shape (increasing utility, decreasing marginal utility) and so it has the standard properties of a CRRA utility function, which makes finding reasonable initial guesses for parameters relatively easy.}

The other main change to this model is the firm side of the economy. We now use a representative firm with a constant-returns-to-scale Cobb-Douglas production function,
\begin{equation}
 Y= A K^{\alpha} L^{1-\alpha}
\end{equation}
where $A$ is aggregate productivity (which will just be normalized to one, we will want it in a future model). $K$ is aggregate physical capital, and is given by $K=\int a d\mu$. Note that $\mu$ is now a function of both the household states $(a,j)$.

We assume there are competitive capital and labor markets, and so in general equilibrium it follows that the interest rate is equal to the marginal product of capital (net of depreciation\footnote{In the budget constraint we have $(1+r)a$, so the interest rate is net of depreciation and therefore equals the marginal product of capital net of depreciation. We could alternatively have used $(1+r-\delta)a$ in the budget constraint in which case the interest rate would be the gross interest rate and equal to the marginal product of capital. $\delta$ is the depreciation rate.}) and the wage is equal to the marginal product of labor, which gives us the general equilibrium condition,
\begin{eqnarray*}
 r=\alpha A K^{\alpha-1} L^{1-\alpha}-\delta \\
 w=(1-\alpha)A K^{\alpha} L^{-\alpha}
\end{eqnarray*}
where still $L= \int \kappa_j h d\mu$.

Because agents are now born with a certain amount of assets we also need to make an assumption about this. We will assume that all agents are born with zero assets, so $\mu(0,1)=1$, and $\mu(a,1)=0$ for all $a \neq 0$. This means we need to change $jequaloneDist$ which represents the distribution of newborns in the codes.

The definition of a stationary equilibrium for this model is now,
\begin{definition}
 A stationary equilibrium is price $r$ (and $w$), and government policies $\tau$ and $pension$ such that
 \begin{enumerate}
  \item Given prices ($r$ and $w$) and government policies ($\tau$, $pension$), the household value function, $V$, and related policy function, $policy$, solve the household problem.
  \item The agents distribution evolves according to,
   \\  $\mu(a_{j+1},j+1)=\int \mathbb{I}_{a_{j+1}=policy_{a'}} \frac{\mu^{j+1}}{\mu^j} d\mu(a_j,j), \; j=1,2,...,J,\text{ and } \mu(:,1)=jequaloneDist$.
  \item The aggregates are based on household policies and agent distribution,
   \\ $L=\int \kappa_j policy_h d\mu$
   \\ $K=\int a d\mu$
   \\ $PensionSpending=\int pension \mathbb{I}_{j\geq Jr} d\mu$
  \item Labor markets clear: $w=(1-\alpha)A K^{\alpha}L^{-\alpha}$
   \\ Capital markets clear: $r=\alpha A K^{\alpha-1} L^{1-\alpha} - \delta$
   \\ Pension system balance: $PensionSpending=\tau*w*L$
   \\ Accidental bequest balance: $AccidentBeq=(\int policy_{a'}(1-s_j) d\mu)/(1+n)$
 \end{enumerate}
\end{definition}
\noindent note the difference between $\mu(a,j)$ which is the population distribution at age $j$, and $\mu^j$ which is the marginal population distribution at age $j$; there is now a meaningful difference. Note how the $policy$ for next period household assets now helps determine the evolution of the agent distribution.

Notice that for accidental bequest balance we have equation $AccidentBeq=(\int policy_{a'} (1-s_j) d\mu)/(1+n)$ which has the lump-sum given to households on the left, being equal to the assets of the dying on the right. The scaling by $1/(1+n)$ is due to the population growth that occurs between last period and this period.\footnote{This formula implicitly assumes $s_J=0$.}

\subsection{OLG Model 5: Progressive Taxation and Government Budget Balance}
We introduce Government in the form of a progressive income tax to raise revenue that is spent on government consumption. As well as including the taxation in the household problem we will also get a new general equilibrium condition relating to the government budget balance. We do not yet have government debt so we simply enforce that the government runs a balanced budget in every period. Introducing government debt is left as an Assignment, see Section \ref{sec:govdebt}.

We will model a progressive income tax, this will involve first defining income, which should be done in accordance with the tax system of the country being modelled. We define $Income=r*a+w*\kappa_j*h$, the sum of capital income and wage earnings, which is close to the definition for income taxes in most countries (note that for retirees $h=0$). We will then use the tax function $IncomeTax=\eta_1+\eta_2*log(Income)*Income$, where $IncomeTax$ is the amount paid by a household with $Income$.\footnote{This functional form is found to have a good empirical fit to the US income tax system by \cite*{GunerKaygusuvVentura2014}. Work on fitting functions for taxation often focuses on 'effective' tax rates, the amount of tax actually paid as a percentage of income, rather than 'statutory' tax rates, what the law says the percentage tax rate is. Effective tax rates have the advantage that the model will reflect how much tax is actually paid, and are a reduced-form way of capturing tax avoidance, e.g., by taking tax-deductions (tax avoidance is used to describe legal ways of paying less tax, illegal ways are called tax evasion). Older work often used the functional form of \cite*{GouveiaStrauss1994,GouveiaStrauss1999}, $IncomeTax = \eta_1 \left[Income-(Income^{-\eta_2}+\eta_3)^{-1/\eta_2}\right]$. Observe that with this functional form $\eta_1$ defines the top (asymptotic) marginal tax rate, while $\eta_2$ and $\eta_3$ control the curvature and initial steepness. In the notation of \cite*{GouveiaStrauss1999} $\eta_1=b$, $\eta_2=p$, $\eta_3=s$. Another functional form that is sometimes used is from \cite*{HeathcoteStoreslettenViolante2014} who use the form: $IncomeTax=Income-\lambda Income^{1-\tau}$, where (1-)$\lambda$ is a parameter controlling the level of taxation and $\tau$ is a parameter controlling the progressivity (this is more normally written as after-tax income equaling $\lambda Income^{1-\tau}$).}

We just model government spending as government consumption, which is indistinguishable from private consumption in this model. There are plenty of alternatives, such as adding a 'public capital' term to the production function, which provide a more explicit purpose to government spending, or to model public sector employment. We here stick with a simple government consumption setting purely for ease of exposition; it allows us to add government without having to make any other substantial changes to our model.\footnote{Note that in the current model there is no 'point' to the Government spending and so the efficient level of government spending would be zero. We can think of it as satisfying some unmodelled purpose, and so we just set it exogenously to some level. The alternative would be to modify the model to create an explicit purpose for government spending, such as some form of social insurance or public capital.}

We need some way to decide 'how big' government should be. We could do this by making sure that the ratio of government revenues to GDP is the same in the model as it is empirically, or to do the same thing but for government spending. But because our model includes just a part of what constitutes government as a whole we won't do this as it would lead to a model with a tiny government. Instead we will do this be deciding how much government spending to have, and then determine taxation to balance the government budget constraint. We do this by setting a target for government consumption as a fraction of GDP, called $GdivYtarget$. We then set $\eta_1$ in general equilibrium to get income tax revenue to be the amount that balances the government budget.\footnote{This will 'sacrifice' the accuracy of the tax rates at the household level, but since our model does not have a realistic income distribution, something we will discuss in a later model, we cannot have both accurate effective tax rates and accurate total tax revenues.}

So our household problem is now modified to include the progressive income tax
\begin{eqnarray*}
 V(a,j)&=& \max_{c,h,a'} \frac{c^{1-\sigma}}{1-\sigma} - \psi \frac{h^{1+\eta}}{1+\eta}  + \beta s_j V(a',j+1) \\
&&\text{if working, }j<=Jr: Income=ra+w \kappa_j h \\
&&\text{if retired, }j\geq Jr: Income=ra \\
&& IncomeTax=\eta_1+\eta_2*log(Income)*Income \\
&&\text{if working, }j<=Jr: c+a'=(1-\tau) w \kappa_j h + (1+r) a -IncomeTax +(1+r)AccidentBeq\\
&&\text{if retired, }j\geq Jr: c+a'=pension+ (1+r) a -IncomeTax +(1+r)AccidentBeq +\mathbb{I}_{j=J} warmglow(a') \\
&& 0\leq h \leq 1
\end{eqnarray*}
Notice that payroll tax is paid on 'pre-income-tax' labor earnings, whether this is appropriate depends on details of the tax system being modelled.

We also need to define the tax revenue as,
\begin{equation*}
 IncomeTaxRevenue=\int (\eta_1+\eta_2*log(Income)*Income) d\mu
\end{equation*}
where $Income$ is as defined above and depends both on the household state (through age (working/retired) and assets) and policies (through hours worked).

The government budget balance is then given by,
\begin{equation*}
 G=IncomeTaxRevenue
\end{equation*}
notice that we are modelling the pensions and the rest of Government as two separate budgets. Whether this is appropriate depends on the country being modelled.

The definition of a stationary equilibrium for this model is now,
\begin{definition}
 A stationary equilibrium is price $r$ (and $w$), and government policies $\tau$ and $pension$ such that
 \begin{enumerate}
  \item Given prices ($r$ and $w$) and government policies ($\tau$, $pension$), the household value function, $V$, and related policy function, $policy$, solve the household problem.
  \item The agents distribution evolves according to,
   \\  $\mu(a_{j+1},j+1)=\int \mathbb{I}_{a_{j+1}=policy_{a'}} \frac{\mu^{j+1}}{\mu^j} d\mu(a_j,j), \; j=1,2,...,J,\text{ and } \mu(:,1)=jequaloneDist$.
  \item The aggregates are based on household policies and agent distribution,
   \\ $L=\int \kappa_j policy_h d\mu$
   \\ $K=\int a d\mu$
   \\ $PensionSpending=\int pension \mathbb{I}_{j\geq Jr} d\mu$
   \\ $IncomeTaxRevenue=\int (\eta_1+\eta_2*log(Income)*Income) d\mu$
  \item Labor markets clear: $w=(1-\alpha)A K^{\alpha}L^{-\alpha}$
   \\ Capital markets clear: $r=\alpha A K^{\alpha-1} L^{1-\alpha} - \delta$
   \\ Pension system balance: $PensionSpending=\tau*w*L$
   \\ Government budget balance: $G=IncomeTaxRevenue$
   \\ Accidental bequest balance: $AccidentBeq=(\int policy_{a'}(1-s_j) d\mu)/(1+n)$
 \end{enumerate}
\end{definition}
\noindent note that we have added an aggregate and a general equilibrium condition, and that the household problem is slightly different.\footnote{When implementing we need to substitute for $Income$ with its expression in terms of the household variables (see the household problem), as otherwise the integral over the agent distribution does not make sense. I do not do so here purely because otherwise the equation gets a bit messy.}

We also need to enforce the condition that government consumption as a fraction of GDP is equal to $GdivYtarget$, and the easiest way is to include this as another general equilibrium condition.\footnote{That $G=GdivYtarget*Y$.} We do this in the codes. 

\subsubsection{Assignment 1: Introduce Government Debt} \label{sec:govdebt}
Assignment 1 is to introduce government debt into OLG Model 5. You will need to create a parameter, call it $B$, which is the amount of government debt (b for bonds). The key insight is that total household assets are no longer equal to aggregate capital, as households will own both capital and the government debt in the form of bonds. So you need to make it so that the aggregate assets held by households is now called something else, e.g. Assets, and modify the FnsToEvaluate appropriately. Notice that then $K=Assets-B$ and you can just replace $K$ with $Assets-B$ in any relevant general equilibrium conditions. You will also need to modify the government budget constraint to include the interest expenses on the debt, so the Government budget balance becomes $G+r*B=IncomeTaxRevenue$.

The assignment is to take the OLG Model 5 codes and modify them to do this. You need to create the parameter for $B$, and make a small modification to the FnsToEvaluate (change from $K$ to $Assets$), and make a few changes to the general equilibrium conditions. You also need to slightly modify the model output that is being reported so that the concepts are what they say they are. (To check how you do a code implementing this is provided in the Assignments folder that you can use to compare to your own version.)










\subsection{OLG Model 6: Stochastic OLG}
We now switch to a OLG with stochastic idiosyncratic shocks. We will introduce a variety of details at once, if you want to understand what they are doing individually check the \href{https://www.vfitoolkit.com/updates-blog/2021/an-introduction-to-life-cycle-models/}{Introduction to Life-Cycle Models}. We won't make any changes to the firm and markets from what it was in OLG Model 5. The ordering of concepts in VFI Toolkit means we think of the state space as being: endogenous state(s), exognenous state(s), model period/age.

We will add both a persistent and a transitory component for (hourly) earnings shocks. The persistent shock is an AR(1) and the transitory shock is i.i.d. normal. This is probably the most standard process in the literature.\footnote{There are plenty of ways to improve on this, but it gives a reasonable baseline to start.} VFI Toolkit calls markov variables, 'z' variables, and also has a feature to handle i.i.d variables, which it calls 'e' variables.\footnote{Alternatively, VFI Toolkit can simply treat an iid variable as if it were markov (by creating a transition matrix for which all the rows are identical). Both approaches give the same answer, but by letting VFI Toolkit know that the shock is i.i.d. the code is able to take advantage of this and run faster.} Papers that empirically estimate earnings processes like those used here often also include a fixed effect, OLG Model 10 demonstrates how to extend the present model to do precisely this.

The household problem to be solved is unchanged except for the introduction of the two shocks.
\begin{eqnarray*}
 V(a,z,e,j)&=& \max_{c,h,a'} \frac{c^{1-\sigma}}{1-\sigma} - \psi \frac{h^{1+\eta}}{1+\eta}  + \beta s_j E_{z',e'}[V(a',j+1,z',e')|z] \\
&&\text{if working, }j<=Jr: Income=ra+w* \kappa_j*exp(z+e) h \\
&&\text{if retired, }j\geq Jr: Income=ra \\
&& IncomeTax=\eta_1+\eta_2*log(Income)*Income \\
&&\text{if working, }j<=Jr: c+a'=(1-\tau) w \kappa_j*exp(z+e) h + (1+r) a -IncomeTax +(1+r)AccidentBeq\\
&&\text{if retired, }j\geq Jr: c+a'=pension+ (1+r) a -IncomeTax +(1+r)AccidentBeq +\mathbb{I}_{j=J} warmglow(a') \\
&& 0\leq h \leq 1 \\
&& z'=\rho_z z + \epsilon, \quad \epsilon \sim N(0,\sigma_{\epsilon,z}^2) \\
&& e \sim N(0,\sigma_{e}^2)
\end{eqnarray*}
Note that we have $\kappa_j*exp(z+e)$, many papers would instead use $exp(\kappa_j+z+e)$ where $\kappa_j$ would now be the log of the values we use here;\footnote{Or you could of course have $\kappa_j*z*e$ by redefining $z$ and $e$ to be the exponentials of what we use here. In the code this just means either taking an exponential or log of the grid, and then altering things like the ReturnFn and FnsToEvaluate appropriately.} I leave it unchanged so that the numbers in the codes can just be left the same. The state vector of the value function now includes both the exogenous shock processes, and next period value function is in expectation. Notice that the expectation is over both $z'$ and $e'$, and that because $z$ is markov it is conditional on this period value of $z$ (whereas since $e$ is i.i.d. the this period value is not relevant to the expectations for the next period realisation).

We also need to define the tax revenue as,
\begin{equation*}
 IncomeTaxRevenue=\int (\eta_1+\eta_2*log(Income)*Income) d\mu
\end{equation*}
where $Income$ is as defined above and depends both on the household state (through age (working/retired) and assets) and policies (through hours worked).

The government budget balance is then given by,
\begin{equation*}
 G=IncomeTaxRevenue
\end{equation*}
notice that we are modelling the pensions and the rest of Government as two separate budgets. Whether this is appropriate depends on the country being modelled.

The definition of a stationary equilibrium for this model is now,
\begin{definition}
 A stationary equilibrium is price $r$ (and $w$), and government policies $\tau$ and $pension$ such that
 \begin{enumerate}
  \item Given prices ($r$ and $w$) and government policies ($\tau$, $pension$), the household value function, $V$, and related policy function, $policy$, solve the household problem.
  \item The agents distribution evolves according to,
   \\  $\mu(a_{j+1},z_{t+1},e_{t+1},j+1)=\int \mathbb{I}_{a_{j+1}=policy_{a'}} \frac{\mu^{j+1}}{\mu^j} d\mu(a,z,e,j) \Gamma(z_{t+1}|z_t) Pr(\epsilon_{t+1}), \; j=1,2,...,J,\text{ and } \mu(:,:,:,1)=jequaloneDist$.
  \item The aggregates are based on household policies and agent distribution,
   \\ $L=\int \kappa_j exp(z+e) policy_h d\mu$
   \\ $K=\int a d\mu$
   \\ $PensionSpending=\int pension \mathbb{I}_{j\geq Jr} d\mu$
   \\ $IncomeTaxRevenue=\int (\eta_1+\eta_2*log(Income)*Income) d\mu$
   \\ $AccidentalBeqLeft= \int policy_{a'}(1-s_j) d\hat{\mu}$
  \item Labor markets clear: $w=(1-\alpha)A K^{\alpha}L^{-\alpha}$
   \\ Capital markets clear: $r=\alpha A K^{\alpha-1} L^{1-\alpha} - \delta$
   \\ Pension system balance: $PensionSpending=\tau*w*L$
   \\ Government budget balance: $G=IncomeTaxRevenue$
   \\ Accidental bequest balance: $AccidentBeq=AccidentalBeqLeft/(1+n)$
 \end{enumerate}
\end{definition}
\noindent note that the household problem is slightly different, and that the agent distribution (and hence all the integrals) is also now over the state space that includes the idiosyncratic exogenous shocks, namely over $(a,z,e,j)$.\footnote{When implementing we need to substitute for $Income$ with its expression in terms of the household variables (see the household problem), as otherwise the integral over the agent distribution does not make sense. I do not do so here purely because otherwise the equation gets a bit messy.} Note that the evolution of the agent distribution has been modified to also include the probabilities of the shocks.

Because we assume a continuum of agents, the idiosyncratic shocks experienced by the indidivual agents do not create any uncertainty for the aggregate economy because we integrate over the entire distribution.\footnote{Intuitively, if you have the whole population then you know all the distribution's moments/statistics with precision.} Hence aggregates like $L$ and prices like $w$ are still just numbers like they were in the model without idiosyncratic shocks.

We also need to enforce the condition that government consumption as a fraction of GDP is equal to $GdivYtarget$, and the easiest way is to include this as another general equilibrium condition. We do this in the codes.

\subsection{Incomplete Markets in General Equilibrium}
If you compare OLG Model 5 and OLG Model 6 you will see that the equilibrium has a higher aggregate capital-output ratio (K/Y) and a lower equilibrium real interest rate. This reflects one of the most important effects of incomplete markets, namely that they lead to precautionary savings and that as a result there is more capital (lower interest rate) than in the complete markets economy. Because the complete markets economy essentially defines efficiency the incomplete markets economy is inefficient because it has too much capital. It is for this reason (among other reasons) that incomplete markets models give results like: that the optimal capital income taxation rate is greater than zero (as it reduces capital towards the efficient level for low tax rates), that the optimal level of government debt is positive (as it soaks up some savings that would otherwise go into capital), that optimal income taxes are progressive (as they provide a form of insurance that reduces the need for precautionary savings), and that some kind of transfers to low asset households is efficient (again reducing the need for precautionary savings). None of these would be the case in a Representative Agent model, but they are for incomplete market heterogeneous agent models.

\subsection{OLG Model 7: Analyzing the OLG model}
We will just reuse OLG model 6, but look at a variety of model outputs, including aggregates, life-cycle profiles, simulating panel data and running a regression on it, and inequality. Since the model is unchanged we will not repeat it here.

Once you have found the general equilibrium we can analyse it using the same kinds of functions that are used for Life-Cycle models. We will simulate some panel data and run a regression on this that measures the amount of 'consumption insurance against earnings shocks' (this model will have less consumption insurance than in the data).

There are dedicated commands for looking at inequality statistics and we can see that while the Gini coefficient of earnings is not too far from realistic (given we haven't really calibrated the model appropriately) the top earnings inequality (such as the share of earnings of the top 5\% or top 1\%) are much smaller than in the data.



\subsection{OLG Model Examples based on papers}
You can find example codes of some OLG models that implement the 'baseline' models of \cite*{ImrohorogluImrohorogluJoines1995}, \cite*{Huggett1996}, \cite*{ConesaKrueger1999}, and \cite*{ConesaKitaoKrueger2009} at:
\\ \href{https://github.com/vfitoolkit/VFItoolkit-matlab-examples/tree/master/OLG}{https://github.com/vfitoolkit/VFItoolkit-matlab-examples/tree/master/OLG}

If you have followed all the OLG models until this point then you will be able to understand both what those papers do and how the codes are going about implementing the baseline models.

\subsection{OLG Model 8: Deterministic Economic/productivity growth}
We now add deterministic productivity growth to the model. To be able to solve it we first have to renormalize the model by 'dividing' everything by the growth to get a stationary model, and then solve this. We can then put the growth back into the solution, should we wish to do so (alternatively, people will often first remove the growth/trend from the empirical data before comparing it to the model). Intuitively, this is just the same trick as when solving the Solow growth model by looking at the balanced growth path (by switching things into per capita per technology unit we get a stationary model, and can then solve this), and in fact the way that we renormalize the model will be exactly analagous. To be able to include growth we have to change the utility function to be non-seperable (that the marginal utility of leisure must depend on the marginal utility of consumption; intuitively, otherwise economic growth leads people to be so rich they never work); for an explanation of this see Life-Cycle Model 22 in the \href{https://www.vfitoolkit.com/updates-blog/2021/an-introduction-to-life-cycle-models/}{\textcolor{blue}{Introduction to Life-Cycle Models}}. We will use almost the same model as OLG Model 6, but with the introduction of deterministic productivity growth, the change to non-seperable utility function, and indexing pensions to growth.

We first write out the model with growth, in which there are now necessarily time subscripts everywhere, and then transform this into the stationary model which we can actually compute and solve (and if desired we can then reinflate the results, we will do this towards the end). We will have deterministic productivity growth at rate $g$, and population growth at rate $n$. We will begin with a model that includes both of these, and then convert to per-capita per-technology units to get a model which has a stationary equilibrium we can solve.

As we will see below when looking at the household problem, before we renormalize the value function $V_t$ now depends on $t$ (and not just $j$) and so we cannot solve for it in the same way as normal. Similarly in the model with growth we cannot define the stationary equilibrium in the same way we did previously. It is still possible to define a general equilibrium in such a model, but this is not very useful to us computationally (Appendix \ref{app:eqmwithgrowth} provides the definition). Instead we need to 'remove' the growth from the model by renormalizing it so that the resulting normalized model has a stationary equilibrium, and we will then be able to solve for this in the same way as usual. The basic idea is that because all of the growth in the economy comes from the $(1+g)^t$ term (in $A_t$ in the production function) we can just divide everything be $(1+g)^t$ and then the resulting model will not have any growth; implementing this idea is slightly more subtle, but that is the concept. Let's do it!

Deterministic productivity growth is introduced using a productivity term in the Cobb-Douglas production function, which is now
\begin{equation*}
 Y_t=A K_t^{\alpha} ((1+g)^t L_t)^{1-\alpha}
\end{equation*}
where $g$ is the growth rate. Notice that instead of putting $A$ out the front we could have put it in with the $(1+g)^t$ term (it would just be $\hat{} \,( 1/(1-\alpha))$ of it's current value).\footnote{We could alternatively have had the $(1+g)^t$ term out front as total-factor productivity growth (rather than labor augmenting technology growth as we have here), but it would then simply complicate the formulas for normalization (and we would no longer get that $g$ is also the growth rate of output-per-capita, which we will here; note that output-per-capita growth is TFP growth together with capital deepening in this model). There is no meaningful difference between TFP growth and labor-augmenting productivity growth when we have a Cobb-Douglas production function.}

To be able to solve this model we need to convert it into per-capita per-technology unit terms. To this end, define the following,
\begin{eqnarray*}
 \hat{Y} &=& \frac{Y}{(1+g)^t (1+n)^t N_0} \\
 \hat{K} &=& \frac{K}{(1+g)^t (1+n)^t N_0} \\
 \hat{L} &=& \frac{L}{(1+n)^t N_0} \\
 \hat{w} &=& \frac{w}{(1+g)^t} \\
\end{eqnarray*}
where the hat, $\hat{}$, denotes per-capita per-technology unit for $Y$ and $K$, but just per capita for $L$, and just per-technology unit for $w$. Notice that as well as having population growth rate $n$ we have some initial population $N_0$ (which previously was implicitly normalized to 1), thus our $\hat{L}$ is labor supply per-capita which is what we had been implicitly using in all of our models up until now.\footnote{It is written here as if I have defined these per-capita per-technology unit terms, and will now use them in the production function. In fact the formulation of these terms follows from how the production function is defined, and the choice to model labor-augmenting technology; if we used something other production function, or if we used, e.g., TFP technology, then we would likely have to use a different formulation to 'remove' the growth from the model. Essentially, when creating the model I first wrote down the production function and then backed out these expressions for how to define the 'hat' variables from this.}

Now that we have these, let's first see how the production function can be rewritten in terms of these variables (even though we do not actually need it to solve the model), and then we will look at how the expressions for the general equilibrium conditions change. We will then deal with the household problem and the evolution of the agent distribution.  We will lastly look at the formulas for the aggregates.\footnote{We need to know both the renormalizations for the aggregates themselves, and the renormalizations for the household-level variables, which we won't know until we have thought about how to handle both the production function and the household problem.}

The production function is $ Y_t=A K_t^{\alpha} ((1+g)^t L_t)^{1-\alpha}$, and we can substitute in this using the (rearranged) definitions of the 'hat' variables,
\begin{eqnarray*}
  Y_t=A K_t^{\alpha} ((1+g)^t L_t)^{1-\alpha} \\
  \hat{Y}_t (1+g)^t (1+n)^t N_0 = A [\hat{K}_t (1+g)^t (1+n)^t N_0 ]^{\alpha} [(1+g)^t \hat{L} (1+n)^t N_0]^{1-\alpha} \\
  \hat{Y}_t (1+g)^t (1+n)^t N_0 = A \hat{K}_t^{\alpha} \hat{L}_t^{1-\alpha} (1+g)^t (1+n)^t N_0 \\
  \hat{Y}_t = A \hat{K}_t^{\alpha} \hat{L}_t^{1-\alpha}
\end{eqnarray*}
Recall that we get expressions for labor market clearance and capital market clearance from the requirement that the wage equals the marginal product of labor, and the interest rate (net of depreciation) equals the marginal product of capital (because we assume perfect competition in labor and capital markets). Taking the derivative of ouput with respect to labor we get the first line of the following, and can then substitute and simplify,
\begin{eqnarray*}
 w_t=(1-\alpha) A K_t^{\alpha}  ((1+g)^t L_t)^{-\alpha} (1+g)^t \\
 \hat{w}_t (1+g)^t = (1-\alpha) A [\hat{K}_t (1+g)^t (1+n)^t N_0 ]^{\alpha} [(1+g)^t \hat{L} (1+n)^t N_0]^{-\alpha} (1+g)^t \\
 \hat{w}_t (1+g)^t = (1-\alpha) A \hat{K}_t^{\alpha} \hat{L}_t^{-\alpha} (1+g)^t \\
 \hat{w}_t = (1-\alpha) A \hat{K}_t^{\alpha} \hat{L}_t^{-\alpha}
\end{eqnarray*}
and doing similarly for the interest rate
\begin{eqnarray*}
 r_t=\alpha A K_t^{\alpha-1}  ((1+g)^t L_t)^{1-\alpha} -\delta \\
 r_t = (1-\alpha) A [\hat{K}_t (1+g)^t (1+n)^t N_0 ]^{\alpha-1} [(1+g)^t \hat{L} (1+n)^t N_0]^{1-\alpha} \\
 r_t = (1-\alpha) A \hat{K}_t^{\alpha-1} \hat{L}_t^{1-\alpha}
\end{eqnarray*}
Notice that we still just have $r_t$. We could define $\hat{r}_t=r_t$ if we wanted to have it so everything in our renormalized model was expressed with a 'hat', but we won't bother to do that here.

We have already dealt with two of the general equilibrium conditions, so now let's handle the other three: pension system balance, government budget balance, and accidental bequests. Pension system balance is trivial if we assume that the household pensions are indexed to per-capita income growth, and so total pension spending growth with $n$ and $g$. We get that pension system balance is now just $PensionSpending=\tau w L_t$, which becomes $\widehat{PensionSpending} (1+g)^t (1+n)^t N_0=\tau \hat{w} (1+g)^t (1+n)^t N_0 \hat{L}_t$, where we define $\widehat{PensionSpending}=PensionSpending/((1+g)^t (1+n)^t N_0)$ in per-capita per-technology unit terms. The Pension system balance simplifies to $\widehat{PensionSpending}=\tau \hat{w} \hat{L}_t$. Similarly we get $\hat{G}=\widehat{IncomeTaxRevenue}$, where we define $\hat{G}=G/((1+g)^t (1+n)^t N_0)$ and  $\widehat{IncomeTaxRevenue}=IncomeTaxRevenue/((1+g)^t (1+n)^t N_0)$. We need to slightly rework the accidental bequests balance, which we now set as $AccidentBeq_t=AccidentBeqLeft_{t-1}$ (previously we used $AccidentBeqLeft/(1+n)$ in the codes, but not in the definition, where the $/(1+n)$ was relflecting that it related to the previous periods and the model was in per-capita terms). We can renormalize this as $\widehat{AccidentBeq}_t=\widehat{AccidentBeqLeft}_{t-1}/((1+g)(1+n))$, where the 'hat' are both defined by dividing by $(1+g)^t (1+n)^t N_0$, note that the $/((1+g)(1+n))$ comes from the different time subscripts.

We are now ready to handle the household problem, which is given by
\begin{eqnarray*}
 V_t(a_t,z,e,j)&=& \max_{c_t,h,a_{t+1}} \frac{c_t^{1-\sigma_1} (1-h_t)^{1-\sigma_1}}{1-\sigma_2} + \beta s_j E_{z',e'}[V_{t+1}(a_{t+1},j+1,z',e')|z] \\
&&\text{if working, }j<=Jr: Income_t=r_t a_t+w_t* \kappa_j*exp(z+e) h \\
&&\text{if retired, }j\geq Jr: Income_t=ra_t \\
&& IncomeTax_t=\eta_1+\eta_2*log(Income_t)*Income_t \\
&&\text{if working, }j<=Jr: c_t+a_{t+1}=(1-\tau) w_t \kappa_j*exp(z+e) h + (1+r) a_t -IncomeTax_t +(1+r)AccidentBeq \\
&&\text{if retired, }j\geq Jr: c_t+a_{t+1}=pension_t+ (1+r_t) a_t -IncomeTax_t +(1+r_t)AccidentBeq_t +\mathbb{I}_{j=J} warmglow(a_{t+1}) \\
&& 0\leq h \leq 1 \\
&& z'=\rho_z z + \epsilon, \quad \epsilon \sim N(0,\sigma_{\epsilon,z}^2) \\
&& e \sim N(0,\sigma_{e}^2)
\end{eqnarray*}
where importantly $V_t$ is now indexed by $t$. I have put $t$ subscipts only on the variables that will grow with time (hence they are on things like wage, assets and consumption, but not on things like idiosyncratic shocks and fraction of time worked).

To see how we are going to renormalize this, it helps to focus on the earnings $w_t* \kappa_j*exp(z+e) h$. Notice how the growth in earnings is going to be captured in $w_t$, and so we don't have any growth in the $\kappa_j*exp(z+e) h$ (effective hours worked) component. We already defined $\hat{w}=\frac{w}{(1+g)^t}$ above when dealing with the production function, and so if we divide the whole budget constraint by $(1+g)^t$ we will be able to replace the $w_t$ with $\hat{w}_t$ and the earnings term will no longer be growing. This intuition is the basis for how we define the rest of our concepts: we already saw that $r_t$ doesn't growth, so for the $(1+r_t)a_t$ term we will define $\hat{a}_t=a_t/((1+g)^t)$. Repeating this approach we get $\hat{c}_t=c_t/((1+g)^t)$, $\widehat{Income}_t=Income_t/((1+g)^t)$, $\widehat{IncomeTax}_t=IncomeTax_t/((1+g)^t)$, $\widehat{AccidentBeq}_t=AccidentBeq_t/((1+g)^t)$,  $\widehat{pension}_t=pension_t/((1+g)^t)$.

Making all of these substitutions we get that the household problem becomes
\begin{eqnarray*}
 V(\hat{a},z,e,j)&=& \max_{\hat{c},h,\hat{a}'} \frac{\hat{c}^{1-\sigma_1} (1-h)^{1-\sigma_1}}{1-\sigma_2} + \beta (1+g)^{\sigma_1(1-\sigma_2)} s_j E_{z',e'}[V(\hat{a}',j+1,z',e')|z] \\
&&\text{if working, }j<=Jr: \widehat{Income}=r \hat{a}+\hat{w}* \kappa_j*exp(z+e) h \\
&&\text{if retired, }j\geq Jr: \widehat{Income}=r\hat{a} \\
&& \widehat{IncomeTax}=\hat{\eta}_1+\eta_2*log(\widehat{Income})*\widehat{Income} \\
&&\text{if working, }j<=Jr: \hat{c}+(1+g)\hat{a}'=(1-\tau) \hat{w} \kappa_j*exp(z+e) h + (1+r) \hat{a} -\widehat{IncomeTax} +(1+r)\widehat{AccidentBeq} \\
&&\text{if retired, }j\geq Jr: \hat{c}+(1+g)\hat{a}'=pension_t+ (1+r_t) \hat{a} -\widehat{IncomeTax} +(1+r)\widehat{AccidentBeq} +\mathbb{I}_{j=J} warmglow((1+g)\hat{a}') \\
&& 0\leq h \leq 1 \\
&& z'=\rho_z z + \epsilon, \quad \epsilon \sim N(0,\sigma_{\epsilon,z}^2) \\
&& e \sim N(0,\sigma_{e}^2)
\end{eqnarray*}
notice most importantly that the discount factor is now $\beta (1+g)^{\sigma_1(1-\sigma_2)}$, to understand why see Life-Cycle model 22 in the \href{https://www.vfitoolkit.com/updates-blog/2021/an-introduction-to-life-cycle-models/}{Introduction to Life-Cycle models}. From the perspective of computing the model solution, the most important thing is that now the value function $V$ does not depend on time, and so we can solve it in the standard fashion. Notice also that the tax parameter $\eta_1$ has had to be changed to $\hat{\eta}_1$ as it does not make sense to have a lump-sum component of tax unless this grows with incomes (so $\eta_1$ is growing, and $\hat{\eta}_1=\eta_1/((1+g)^t)$ is constant).

We can now consider the agents distribution. I won't go into the details here but essentially instead of writing it as $\mu(a,z,e,j)$ we can just write it as $\hat{\mu}(\hat{a},z,e,j)=\mu(a/((1+g)^t),z,e,j)$ and there is therefore no change that we need to make (other than putting hats on lots of terms) in the expression for the evolution of the agents distribution. Notice that having solved for $\hat{\mu}(\hat{a},z,e,j)$ it is trivial to recover $\mu(a,z,e,j)$ should we wish to do so. In the discussion of the agent distribution in this paragraph I have, for convenience, ignored writing the time subcripts that should appear on $\mu$ and $a$.

We are just left with the expression for the aggregates, it trivally follows that these are given by,
\begin{eqnarray*}
 \hat{L}=\int \kappa_j exp(z+e) policy_h d\hat{\mu} \\
 \hat{K}=\int a d\hat{\mu} \\
 \widehat{PensionSpending}=\int \hat{pension} \mathbb{I}_{j\geq Jr} d\hat{\mu} \\
 \widehat{IncomeTaxRevenue}=\int (\hat{eta}_1 + \eta_2 log(\widehat{Income}) \widehat{Income} ) d\hat{\mu}
\end{eqnarray*}
Note that what I do with the progressive tax function is a bit of an abritrary hack.\footnote{As discussed in OLG Model 5 there are three main functional forms used to model progressive taxation: \cite*{GunerKaygusuvVentura2014}, \cite*{GouveiaStrauss1994,GouveiaStrauss1999} and \cite*{HeathcoteStoreslettenViolante2014}. But all would mean taxes increase as incomes increase. Renormalizing income as an input, as I do here, can help 'solve' this for all three, but is arbitrary. I am not aware of any evidence on what constitutes a 'good' way to extend any of these to accomodate income growth.}

Whew! We are just about done. Now we can just write the definition for stationary general equilibrium and then get to solving :)

The definition of a stationary equilibrium for this model is now,
\begin{definition}
 A stationary equilibrium is price $r$ (and $\hat{w}$), and government policies $\tau$ and $\widehat{pension}$ such that
 \begin{enumerate}
  \item Given prices ($r$ and $\hat{w}$) and government policies ($\tau$, $\widehat{pension}$), the household value function, $V$, and related policy function, $policy$, solve the household problem.
  \item The agents distribution evolves according to,
   \\  $\hat{\mu}(\hat{a}_{j+1},z_{t+1},e_{t+1},j+1)=\int \mathbb{I}_{\hat{a}_{j+1}=policy_{\hat{a}'}} \frac{\hat{\mu}^{j+1}}{\hat{\mu}^j} d\hat{\mu}(\hat{a},z,e,j) \Gamma(z_{t+1}|z_t) Pr(\epsilon_{t+1}), \; j=1,2,...,J,\text{ and } \hat{\mu}(:,:,:,1)=jequaloneDist$.
  \item The aggregates are based on household policies and agent distribution,
   \\ $\hat{L}=\int \kappa_j exp(z+e) policy_h d\hat{\mu}$
   \\ $\hat{K}=\int \hat{a} d\hat{\mu}$
   \\ $\widehat{PensionSpending}=\int \hat{pension} \mathbb{I}_{j\geq Jr} d\hat{\mu}$
   \\ $\widehat{IncomeTaxRevenue}=\int (\hat{\eta}_1+\eta_2*log(\widehat{Income})*\widehat{Income}) d\hat{\mu}$
   \\ $\widehat{AccidentalBeqLeft}= \int policy_{\hat{a}'}(1-s_j) d\hat{\mu}$
  \item Labor markets clear: $\hat{w}=(1-\alpha)A \hat{K}^{\alpha}\hat{L}^{-\alpha}$
   \\ Capital markets clear: $r=\alpha A \hat{K}^{\alpha-1} \hat{L}^{1-\alpha} - \delta$
   \\ Pension system balance: $\widehat{PensionSpending}=\tau*\hat{w}*\hat{L}$
   \\ Government budget balance: $\hat{G}=\widehat{IncomeTaxRevenue}$
   \\ Accidental bequest balance: $\widehat{AccidentBeq}=\widehat{AccidentalBeqLeft}/((1+g)(1+n))$
 \end{enumerate}
\end{definition}
Notice that this is largely just the same definition as we had for OLG Model 6 (the model to which we added deterministic labor-augmenting productivity growth) except that many of the variables now have 'hat's on them, and some other things like the presence of (1+g) in the budget constraint and the change in the discount factor in the household problem.

We can therefore just solve for this stationary general equilibrium using all the same kinds of code (but making the small changes to discount factor etc.) and get the solution for all of the 'hat' variables.

If we want to recover variables like $Y$ or $w$ we can just use the expressions like $\hat{Y} =\frac{Y}{(1+g)^t (1+n)^t N_0}$ and $\hat{w} = \frac{w}{(1+g)^t}$ that we used to define the hat variables (just use them in reverse). When doing so what we set for things like $N_0$ and $A$, which define the 'starting point' are largely arbitrary (we could pick a specific year in the data to target if we wanted, but if this is our goal we should probably be solving for the general equilibrium transition paths, something not covered in this Intro to OLG).

The codes do just this, creating time-series for $Y$ and $w$ (among others) and then plotting them. To save substantial rewriting of the codes I have not attempted to include the 'hat' in the variable names. Instead at the end, when recovering things like $L$ and $w$ I call them 'actual L' and 'actual w'.


\subsection{OLG Model 9: Exotic preferences}
We now resolve OLG Model 6 twice, once modifying the model to use Quasi-hyperbolic discounting in the household problem, and once modifying the model to use Epstein-Zin preferences in the household problem. Both of these will require changing just a few lines of code. Since the only change to the model is in the definition of the household problem, we will not repeat things like the firm problem, the general equilibrium conditions, and the definition of stationary general equilibrium here, see OLG Model 6 for these. For a more complete explanation of these exotic preferences \textcolor{blue}{\href{https://www.vfitoolkit.com/updates-blog/2021/exotic-preferences-epstein-zin-quasi-hyperbolic/}{see this link}}.

Quasi-hyperbolic discounting is a standard way to model impatience. Epstein-Zin preferences separate the 'intertemporal elasticity of substitution' from the 'risk aversion', both of which are determined by the same single parameter using the standard vonNeumann-Morgenstern preferences that we have been using without explicitly mentioning them until now.

\subsubsection{OLG Model 9A: Quasi-hyperbolic discounting}
There are two variations of Quasi-hyperbolic discounting, naive and sophisticated. In naive quasi-hyperbolic discounting the agent 'naively' believes that their future self will not suffer from the same impatience as their present self, that is, they belive their future self will be an exponential discounter (their future self will actually just be a quasi-hyperbolic discounter like themselves. A sophisticated quasi-hyperbolic discounter knows that their future self will also behave as a quasi-hyperbolic discounter and takes this into account.

Solving naive quasi-hyperbolic is simpler (conceptually, in codes they are equally easy). The future self solves the standard household problem of an exponential discounter that we already saw in OLG Model 6, and the value function for which we denote $V$. The naive quasi-hyperbolic discounter thus faces the household problem,
\begin{eqnarray*}
 \tilde{V}(a,z,e,j)&=& \max_{c,h,a'} \frac{c^{1-\sigma}}{1-\sigma} - \psi \frac{h^{1+\eta}}{1+\eta}  + \beta_0 \beta s_j E_{z',e'}[V(a',j+1,z',e')|z] \\
&&\text{if working, }j<=Jr: Income=ra+w* \kappa_j*exp(z+e) h \\
&&\text{if retired, }j\geq Jr: Income=ra \\
&& IncomeTax=\eta_1+\eta_2*log(Income)*Income \\
&&\text{if working, }j<=Jr: c+a'=(1-\tau) w \kappa_j*exp(z+e) h + (1+r) a -IncomeTax +(1+r)AccidentBeq\\
&&\text{if retired, }j\geq Jr: c+a'=pension+ (1+r) a -IncomeTax +(1+r)AccidentBeq +\mathbb{I}_{j=J} warmglow(a') \\
&& 0\leq h \leq 1 \\
&& z'=\rho_z z + \epsilon, \quad \epsilon \sim N(0,\sigma_{\epsilon,z}^2) \\
&& e \sim N(0,\sigma_{e}^2)
\end{eqnarray*}
Notice that there is the additional discount factor $\beta_0$, and that the left-hand side is now $\tilde{V}$, the value function of the naive quasi-hyperbolic discounter, while the right-hand side is the value function of the exponential discounter that they (incorrectly) believe their future self to be.

The sophisticated quasi-hyperbolic discounter is a little more subtle to write out. They know their future self will act as a quasi-hyperbolic discounter (so use the policy function of the quasi-hyperbolic discount problem), but they discount between any two consecutive future periods using the discount factor $\beta$. So we essentially need, at each time horizon, to first get the policy using the quasi-hyperbolic discounting, and then evaluate the value function using the exponential discount factor.

The value function of the sophisticated quasi-hyperbolic discounter, $\hat{V}$, is given by
\begin{eqnarray*}
 \hat{V}(a,z,e,j)&=& \max_{c,h,a'} \frac{c^{1-\sigma}}{1-\sigma} - \psi \frac{h^{1+\eta}}{1+\eta}  + \beta_0 \beta s_j E_{z',e'}[\underbar{V}(a',j+1,z',e')|z] \\
&&\text{if working, }j<=Jr: Income=ra+w* \kappa_j*exp(z+e) h \\
&&\text{if retired, }j\geq Jr: Income=ra \\
&& IncomeTax=\eta_1+\eta_2*log(Income)*Income \\
&&\text{if working, }j<=Jr: c+a'=(1-\tau) w \kappa_j*exp(z+e) h + (1+r) a -IncomeTax +(1+r)AccidentBeq\\
&&\text{if retired, }j\geq Jr: c+a'=pension+ (1+r) a -IncomeTax +(1+r)AccidentBeq +\mathbb{I}_{j=J} warmglow(a') \\
&& 0\leq h \leq 1 \\
&& z'=\rho_z z + \epsilon, \quad \epsilon \sim N(0,\sigma_{\epsilon,z}^2) \\
&& e \sim N(0,\sigma_{e}^2)
\end{eqnarray*}
and we denote the optimal policies that are the argmax of this problem by $\hat{c}$, $\hat{h}$, and $\hat{a}'$. Using these optimal polices we can define,
\begin{eqnarray*}
 \underbar{V}(a,z,e,j)&=& \frac{\hat{c}^{1-\sigma}}{1-\sigma} - \psi \frac{\hat{h}^{1+\eta}}{1+\eta}  + \beta_0 \beta s_j E_{z',e'}[\underbar{V}(\hat{a}',j+1,z',e')|z] \\
&&\text{if working, }j<=Jr: Income=ra+w* \kappa_j*exp(z+e) \hat{h} \\
&&\text{if retired, }j\geq Jr: Income=ra \\
&& IncomeTax=\eta_1+\eta_2*log(Income)*Income \\
&&\text{if working, }j<=Jr: \hat{c}+\hat{a}'=(1-\tau) w \kappa_j*exp(z+e) \hat{h} + (1+r) a -IncomeTax +(1+r)AccidentBeq\\
&&\text{if retired, }j\geq Jr: \hat{c}+\hat{a}'=pension+ (1+r) a -IncomeTax +(1+r)AccidentBeq +\mathbb{I}_{j=J} warmglow(\hat{a}') \\
&& z'=\rho_z z + \epsilon, \quad \epsilon \sim N(0,\sigma_{\epsilon,z}^2) \\
&& e \sim N(0,\sigma_{e}^2)
\end{eqnarray*}
notice that we just evaluate $\underbar{V}$ using the optimal policies, and so there is no maximization involved here.


\subsubsection{OLG Model 9B: Epstein-Zin preferences}
The household problem to be solved is unchanged except for the use of Epstein-Zin preferences.
\begin{eqnarray*}
 \tilde{V}(a,z,e,j)&=& \max_{c,h,a'} \frac{c^{1-\sigma}}{1-\sigma} - \psi \frac{h^{1+\eta}}{1+\eta}  + \beta s_j E_{z',e'}[V(a',j+1,z',e')^{1+\phi}|z]^{\frac{1}{1+\phi}} \\
&&\text{if working, }j<=Jr: Income=ra+w* \kappa_j*exp(z+e) h \\
&&\text{if retired, }j\geq Jr: Income=ra \\
&& IncomeTax=\eta_1+\eta_2*log(Income)*Income \\
&&\text{if working, }j<=Jr: c+a'=(1-\tau) w \kappa_j*exp(z+e) h + (1+r) a -IncomeTax +(1+r)AccidentBeq\\
&&\text{if retired, }j\geq Jr: c+a'=pension+ (1+r) a -IncomeTax +(1+r)AccidentBeq +\mathbb{I}_{j=J} warmglow(a') \\
&& 0\leq h \leq 1 \\
&& z'=\rho_z z + \epsilon, \quad \epsilon \sim N(0,\sigma_{\epsilon,z}^2) \\
&& e \sim N(0,\sigma_{e}^2)
\end{eqnarray*}
where $\phi$ controls the amount of 'additional' risk aversion.

We have to be careful about how warm-glow of bequests is handled when using Epstein-Zin model. See the Intro to Life-Cycle models for an explanation of Epstein-Zin preferences and the options when implementing them.


\subsection{OLG Model 10: Permanent Types 1, Fixed effects}
This model extends OLG Model 6. When earnings processes are estimated from the data it is standard to include a fixed-effect. We now introduce permanent types as the way for VFI Toolkit to model this: permanent type includes but is much more general than fixed type. The toolkit uses $i$ to refer to permanent types and this can be done either as just giving parameters as a vector as in this example, or by giving each permanent type a name, which is done in the next model.

The household problem now includes a fixed-effect, $\gamma_i$, which depends on the permanent type of the agent.
\begin{eqnarray*}
 V_i(a,z,e,j)&=& \max_{c,h,a'} \frac{c^{1-\sigma}}{1-\sigma} - \psi \frac{h^{1+\eta}}{1+\eta}  + \beta s_j E_{z',e'}[V_i(a',j+1,z',e')|z] \\
&&\text{if working, }j<=Jr: Income=ra+w* \kappa_j*exp(\gamma_i+z+e) h \\
&&\text{if retired, }j\geq Jr: Income=ra \\
&& IncomeTax=\eta_1+\eta_2*log(Income)*Income \\
&&\text{if working, }j<=Jr: c+a'=(1-\tau) w \kappa_j*exp(\gamma_i+z+e) h + (1+r) a -IncomeTax +(1+r)AccidentBeq\\
&&\text{if retired, }j\geq Jr: c+a'=pension+ (1+r) a -IncomeTax +(1+r)AccidentBeq +\mathbb{I}_{j=J} warmglow(a') \\
&& 0\leq h \leq 1 \\
&& z'=\rho_z z + \epsilon, \quad \epsilon \sim N(0,\sigma_{\epsilon,z}^2) \\
&& e \sim N(0,\sigma_{e}^2)
\end{eqnarray*}
notice that there are now $i$ separate household problems, and hence we add a $i$ subscript to $V$.

Solving this model is easy enough. First we say how many different permanent types of agents we want, we will use three ($N\_{}i=3$), and we need to create the parameter $\gamma_i$, which is vector of size 3-by-1 (or 1-by-3, doesn't matter). We also need to declare how many of agents of each type there are, or more accurately the distribution of agents over these types. This is a vector of three fractions that add to one, which we will call 'gamma$\_{}$dist', and we need to put the name of this into PTypeDistParamNames.

We also need to change all the command names, which are now the same as before but with $\_{}PType$ at the end (what this means will be obvious in the codes). That is all.

By default the permanent types will be assigned the names ptype001, ptype002, etc. And everything like value functions, agent distributions, and model statistics use these names (so, e.g., V.ptype001 is the value function for the first permanent type).

While a lot changes behind the scenes, realistically all that is involved is just solving everything $N\_{}i$ times, and things like run times and memory usage will reflect this. By default all model statistics will report both the aggregate for the whole economy exactly as before, e.g., AllStats.K.Mean, and also will report all the same statistics conditional on permanent type of the agent, the codes show this for the life-cycle profiles.

We now write out the definition of stationary general equilibrium in the model, notice that we essentially just have the household and agent distribution problems once for each agent type, while the aggregates are now defined over all the agents.
\begin{definition}
 A stationary equilibrium is price $r$ (and $w$), and government policies $\tau$ and $pension$ such that
 \begin{enumerate}
  \item Given prices ($r$ and $w$) and government policies ($\tau$, $pension$), the household value functions, $V_i$, and related policy function, $policy_i$, solve the household problem; for each $i=1,...,N\_{}i$.
  \item The agents distribution evolves according to,
   \\  $\mu_i(a_{j+1},z_{t+1},e_{t+1},j+1)=\int \mathbb{I}_{a_{j+1}=policy_{i,a'}} \frac{\mu_i^{j+1}}{\mu_i^j} d\mu_i(a,z,e,j) \Gamma(z_{t+1}|z_t) Pr(\epsilon_{t+1}), \; j=1,2,...,J,\text{ and } \mu_i(:,:,:,1)=jequaloneDist \; \text{ for all} i=1,2,...,N\_{}i$
  \item The aggregates are based on household policies and agent distribution,
   \\ $L=\sum_i \int \kappa_j exp(\gamma_i +z+e) policy_{h,i} d\mu_i$
   \\ $K=\sum_i \int a d\mu_i$
   \\ $PensionSpending=\sum_i \int pension \mathbb{I}_{j\geq Jr} d\mu_i$
   \\ $IncomeTaxRevenue=\sum_i \int (\eta_1+\eta_2*log(Income)*Income) d\mu_i$
   \\ $AccidentalBeqLeft= \sum_i \int policy_{\hat{a}'}(1-s_j) d\mu_i$
  \item Labor markets clear: $w=(1-\alpha)A K^{\alpha}L^{-\alpha}$
   \\ Capital markets clear: $r=\alpha A K^{\alpha-1} L^{1-\alpha} - \delta$
   \\ Pension system balance: $PensionSpending=\tau*w*L$
   \\ Government budget balance: $G=IncomeTaxRevenue$
   \\ Accidental bequest balance: $AccidentBeq=AccidentalBeqLeft/(1+n)$
 \end{enumerate}
\end{definition}
Note that you could write this more succinctly by putting $i$ as another dimension of the value function, and the agent distribution. We have not done so because the current way of writing the problem (with separate value functions and agent distributions for each $i$, and with aggregates written as a sum over $i$) is closer to what is being done internally and so hopefully helps make clearer what is happening.

The firm side of the problem and the general equilibrium conditions are all unchanged (although of course the aggregates that go into them as inputs are different).




\subsection{OLG Model 11: Permanent Types 2, Male and female households}
This example models two kinds of agents, male and female. Both agents solve the same problem but face different earnings, both in terms of their deterministic life-cycle profile of earnings and in terms of the earnings shock processes they face. Again we are going to adapt OLG Model 6. We will write out the household problem expicitly here, and the rest of the stationary general equilibrium definition is just the same as it was in OLG Model 10 (except of course that there are now 2 permananent agent types, not 3).

The household problem is,
\begin{eqnarray*}
 V_i(a,z,e,j)&=& \max_{c,h,a'} \frac{c^{1-\sigma}}{1-\sigma} - \psi \frac{h^{1+\eta}}{1+\eta}  + \beta s_j E_{z',e'}[V_i(a',j+1,z',e')|z] \\
&&\text{if working, }j<=Jr: Income=ra+w* \kappa_{i,j}*exp(\gamma+z_i+e_i) h \\
&&\text{if retired, }j\geq Jr: Income=ra \\
&& IncomeTax=\eta_1+\eta_2*log(Income)*Income \\
&&\text{if working, }j<=Jr: c+a'=(1-\tau) w \kappa_{i,j}*exp(\gamma+z_i+e_i) h + (1+r) a -IncomeTax +(1+r)AccidentBeq\\
&&\text{if retired, }j\geq Jr: c+a'=pension+ (1+r) a -IncomeTax +(1+r)AccidentBeq +\mathbb{I}_{j=J} warmglow(a') \\
&& 0\leq h \leq 1 \\
&& z_i'=\rho_{z,i} z_i + \epsilon, \quad \epsilon \sim N(0,\sigma_{\epsilon,z,i}^2) \\
&& e \sim N(0,\sigma_{e,i}^2)
\end{eqnarray*}
notice that there are now $i$ separate household problems, in our case two which we will give the names 'male' and 'female'. The differences between male and female households are just about the labor productivity parameters, $\kappa_{i,j}$, $\gamma_i$, $\rho_{z,i}$, $\sigma_{\epsilon,z,i}$, and $\sigma_{e,i}$. Because the parameters for the exogenous shocks differ the shocks themselves will and so I have written $z_i$ and $e_i$ to emphasize this (although I could just drop this notation).

Rather than just specifying the number of permanent types as in OLG Model 10. Instead we will given them names by setting $Names\_i=\{'male','female'\}$. We can then use this names to specify where the two permanent types differ. For example we set $Params.gamma\_i.male=0.1$ and $Params.gamma\_i.female=0.1$. Setting up separate grids is done in much the same way (e.g. $vfoptions.z\_grid\_J.male$). Because we specify names for the permanent types the outputs will use the same names (e.g., $V.male$ and $V.female$ are the value functions for male and female households).


\subsection{OLG Model 12: Married couple households}
This example models just one kind of household, a married couple. The married couple faces four earning shocks (a persistent shock and a transitory shock for each spouse) and makes two labor supply decisions (one for each spouse).

We solve a model of a household in which there are two people (e.g., a married household). They make a joint decision about how much each will work. This involves two decisions variables for the two labor supply choices. We will set up each household member to have effective hours shocks that are a combination of one AR(1) persistent shock and one i.i.d. shock, for each person (so two of each for the household). We will also allow for correlation between the persistent shocks of the two household members, and for correlation between the transitory shocks of the two household members. We will also allow different deterministic age-dependent labor efficiency units, $\kappa_j$, for each spouse.

Our household value function now has $z1$ and $e1$ for the first spouse, and $z2$ and $e2$ for the second spouse. We also have $\kappa_{j,1}$ and  $\kappa_{j,2}$. We add a disutility term for each spouses labor supply, but the household has joint consumption.
\begin{eqnarray*}
 V(a,z_1,z_2,e_1,e_2,j)&=& \max_{c,aprime,h_1,h_2} \frac{c^{1-\sigma}}{1-\sigma} - \psi \frac{h_1^{1+\eta}}{1+\eta}- \psi \frac{h_2^{1+\eta}}{1+\eta} + (1-s_j)\beta  \mathbb{I}_{(j>=Jr+10)} warmglow(aprime) \\
 && \quad \quad \quad \quad+ s_j \beta E[V(aprime,z_1prime,z_2prime,e_1prime,e_2prime,j+1)|z_1,z_2] \\
&& \text{if }j<Jr: \; c+aprime=(1+r)a+w \kappa_{j,1}  z_1 e_1 h_1+w \kappa_{j,2}  z_2 e_2 h_2 \\
&& \text{if }j>=Jr: \; c+aprime=(1+r)a +pension\\
&& 0\leq h \leq 1, aprime\geq 0 \\
&& log([z_1prime; z_2prime])=\rho_{z} log([z_1; z_2]) + \epsilon, \; \epsilon \sim N(0,\Sigma_{\epsilon,z}^2) \\
&& log([e_1,e_2) \sim N(0,\Sigma_{e}^2)
\end{eqnarray*}
Notice that $[z_1,z_2]$ is VAR(1) (in logs) and $[e_1,e_2]$ is i.i.d. normal (in logs).

We are allowing for $z_1$ and $z_2$ to follow a VAR(1) (so both $\rho_z$ and $\Sigma_{\epsilon,z}$ are 2-by-2 matrices). We do not require that $\Sigma_{\epsilon,z}$ be diagonal, so the innovations themselves can be correlated. We similarly allow for $e_1$ and $e_2$ to be correlated (so $\Sigma_{e}$ is not diagonal).

This model demonstrates how to use two decision variables: you need to put them as the first two entries to the ReturnFn and also to all FnsToEvaluate. The model output includes calculating the life-cycle labor supply for each spouse separately, as well as the total household labor supply.

A common variation of the model used in practice would add a fixed cost of working (e.g., $-\mathbb{I}_{h_2>0}*fc$, where $fc$ is a constant) to capture that the second spouse in many households will sometimes supply zero labor.\footnote{Because we have $0\leq h_2 \leq 1$ together with the shape of the utility function it would never be optimal to set $h_2=0$ in the absence of a fixed cost of working non-zero hours.} Note that implementing this is just a simple modification of the ReturnFn.

We will assume that aggregate effective labor is just the sum of the effective labor units of both spouses.

We now write out the definition of stationary general equilibrium in the model.
\begin{definition}
 A stationary equilibrium is price $r$ (and $w$), and government policies $\tau$ and $pension$ such that
 \begin{enumerate}
  \item Given prices ($r$ and $w$) and government policies ($\tau$, $pension$), the household value functions, $V$, and related policy function, $policy$, solve the household problem.
  \item The agents distribution evolves according to,
   \\  $\mu(a_{j+1},z_{1,t+1},z_{2,t+1},e_{1,t+1},e_{2,t+1},j+1)=\int \mathbb{I}_{a_{j+1}=policy_{a'}} \frac{\mu^{j+1}}{\mu^j} d\mu(a,z_1,z_2,e_1,e_2,j) \Gamma(z_{1,t+1},z_{2,t+1}|z_{1,t},z_{2,t}) Pr(\epsilon_{1,t+1},\epsilon_{2,t+1}), \; j=1,2,...,J,\text{ and } \mu(:,:,:,:,:,1)=jequaloneDist$
  \item The aggregates are based on household policies and agent distribution,
   \\ $L=\int \kappa_{j,1} exp(\gamma_1+z_1+e_1) policy_{h_1} + \kappa_{j,2} exp(\gamma_2+z_2+e_2) policy_{h_2} d\mu$
   \\ $K=\int a d\mu$
   \\ $PensionSpending=\int pension \mathbb{I}_{j\geq Jr} d\mu$
   \\ $IncomeTaxRevenue=\int (\eta_1+\eta_2*log(Income)*Income) d\mu$
   \\ $AccidentalBeqLeft= \int policy_{\hat{a}'}(1-s_j) d\mu$
  \item Labor markets clear: $w=(1-\alpha)A K^{\alpha}L^{-\alpha}$
   \\ Capital markets clear: $r=\alpha A K^{\alpha-1} L^{1-\alpha} - \delta$
   \\ Pension system balance: $PensionSpending=\tau*w*L$
   \\ Government budget balance: $G=IncomeTaxRevenue$
   \\ Accidental bequest balance: $AccidentBeq=AccidentalBeqLeft/(1+n)$
 \end{enumerate}
\end{definition}
Where the main difference to earlier examples (such as OLG Model 10) is that the household problem is now more complicated, and the definition of the aggregate (effective) labor $L$.

The firm side of the problem and the general equilibrium conditions are all unchanged (although of course the aggregates that go into them as inputs are different).

Obviously there are some limitations to modelling a married couple in this way. One if that they have just one 'age' between the two spouses.

Because the 'size' of the household problem for the married couple is quite large we need to use settings that are slower but require less memory. Hence we set vfoptions.lowmemory=1 and simoptions.parallel=4.\footnote{Code will loop over e, instead of paralleling over e, when solving value function iteration problem, and then use 'sparse matrix on cpu' when doing agent distribution. The point of VFI Toolkit is you don't really need to understand the details of what these mean nor how these work to be able to take advantage of them :)}

\subsection{OLG Model 13: Married Couples, Single Males and Single Females}
We now consider an economy with three types of household: married couples, single males and single females. The household problem of married couples is exactly as in OLG Model 11, and the household problems of single males and single females are exactly as in OLG Model 12.

So we have the household problem of the married couple given by,
\begin{eqnarray*}
 V_m(a,z_1,z_2,e_1,e_2,j)&=& \max_{c,aprime,h_1,h_2} \frac{c^{1-\sigma}}{1-\sigma} - \psi \frac{h_1^{1+\eta}}{1+\eta}- \psi \frac{h_2^{1+\eta}}{1+\eta} + (1-s_j)\beta  \mathbb{I}_{(j>=Jr+10)} warmglow(aprime) \\
 && \quad \quad \quad \quad+ s_j \beta E[V_m(aprime,z_1prime,z_2prime,e_1prime,e_2prime,j+1)|z_1,z_2] \\
&& \text{if }j<Jr: \; c+aprime=(1+r)a+w \kappa_{j,1}  z_1 e_1 h_1+w \kappa_{j,2}  z_2 e_2 h_2 \\
&& \text{if }j>=Jr: \; c+aprime=(1+r)a +pension\\
&& 0\leq h \leq 1, aprime\geq 0 \\
&& log([z_1prime; z_2prime])=\rho_{z} log([z_1; z_2]) + \epsilon, \; \epsilon \sim N(0,\Sigma_{\epsilon,z}^2) \\
&& log([e_1,e_2) \sim N(0,\Sigma_{e}^2)
\end{eqnarray*}
Notice that $[z_1,z_2]$ is VAR(1) (in logs) and $[e_1,e_2]$ is i.i.d. normal (in logs). Note that a 'm' subscript has been added to the value function to emphasize that this is for the married couple.

We also have two further household problems for the single male and single female households given by,
\begin{eqnarray*}
 V_i(a,z,e,j)&=& \max_{c,h,a'} \frac{c^{1-\sigma}}{1-\sigma} - \psi \frac{h^{1+\eta}}{1+\eta}  + \beta s_j E_{z',e'}[V_i(a',j+1,z',e')|z] \\
&&\text{if working, }j<=Jr: Income=ra+w* \kappa_{i,j}*exp(\gamma_i+z_i+e_i) h \\
&&\text{if retired, }j\geq Jr: Income=ra \\
&& IncomeTax=\eta_1+\eta_2*log(Income)*Income \\
&&\text{if working, }j<=Jr: c+a'=(1-\tau) w \kappa_{i,j}*exp(\gamma_i+z_i+e_i) h + (1+r) a -IncomeTax +(1+r)AccidentBeq\\
&&\text{if retired, }j\geq Jr: c+a'=pension+ (1+r) a -IncomeTax +(1+r)AccidentBeq +\mathbb{I}_{j=J} warmglow(a') \\
&& 0\leq h \leq 1 \\
&& z_i'=\rho_{z,i} z_i + \epsilon, \quad \epsilon \sim N(0,\sigma_{\epsilon,z,i}^2) \\
&& e \sim N(0,\sigma_{e,i}^2)
\end{eqnarray*}
where $i$ is 'male' or 'female'.

The definition of the stationary equilibrium is essentially just a combination of that from OLG Models 11 and 12 (the sum across permanent types from 11, with the sum across labor supply of both genders from 12). I treat the two single household types as permanent type 1 and 2, and the married couple type as permanent type 3 in this definition.
\begin{definition}
 A stationary equilibrium is price $r$ (and $w$), and government policies $\tau$ and $pension$ such that
 \begin{enumerate}
  \item Given prices ($r$ and $w$) and government policies ($\tau$, $pension$), the household value functions, $V_i$, and related policy function, $policy_i$, solve the household problem; for each $i=1,...,N\_{}i$.
  \item The agents distribution evolves according to,
   \\  $\mu_i(a_{j+1},z_{t+1},e_{t+1},j+1)=\int \mathbb{I}_{a_{j+1}=policy_{i,a'}} \frac{\mu_i^{j+1}}{\mu_i^j} d\mu_i(a,z,e,j) \Gamma(z_{t+1}|z_t) Pr(\epsilon_{t+1}), \; j=1,2,...,J,\text{ and } \mu_i(:,:,:,1)=jequaloneDist \; \text{ for all} i=1,2$
   \\  $\mu_i(a_{j+1},z_{1,t+1},z_{2,t+1},e_{1,t+1},e_{2,t+1},j+1)=\int \mathbb{I}_{a_{j+1}=policy_{i,a'}} \frac{\mu_i^{j+1}}{\mu_i^j} d\mu_i(a,z_1,e_1,z_2,e_2,j) \Gamma(z_{1,t+1},z_{2,t+1}|z_{1,t},z_{2,t}) Pr(\epsilon_{1,t+1},\epsilon_{2,t+1})), \; j=1,2,...,J,\text{ and } \mu_i(:,:,:,:,:,1)=jequaloneDist \; \text{ for all} i=3$
  \item The aggregates are based on household policies and agent distribution,
   \\ $L=\sum_{i=1}^2 \int \kappa_j exp(\gamma_i +z+e) policy_{h,i} d\mu_i + \int [\kappa_{j,1} exp(\gamma_1 +z_1+e_1) policy_{h,i,1} + \kappa_{j,2} exp(\gamma_2 +z_2+e_2) policy_{h,i,2}] d\mu_3$
   \\ $K=\sum_i \int a d\mu_i$
   \\ $PensionSpending=\sum_i \int pension \mathbb{I}_{j\geq Jr} d\mu_i$
   \\ $IncomeTaxRevenue=\sum_i \int (\eta_1+\eta_2*log(Income)*Income) d\mu_i$
   \\ $AccidentalBeqLeft= \sum_i \int policy_{\hat{a}'}(1-s_j) d\mu_i$
  \item Labor markets clear: $w=(1-\alpha)A K^{\alpha}L^{-\alpha}$
   \\ Capital markets clear: $r=\alpha A K^{\alpha-1} L^{1-\alpha} - \delta$
   \\ Pension system balance: $PensionSpending=\tau*w*L$
   \\ Government budget balance: $G=IncomeTaxRevenue$
   \\ Accidental bequest balance: $AccidentBeq=AccidentalBeqLeft/(1+n)$
 \end{enumerate}
\end{definition}
Here we have defined the equilibrium using $N\_i$ which is the three types of household. In the codes we will use $Names\_i$ to make it much easier to keep track of them.

For more on the importance of distinguishing single and married households for macroeconomic aggregates, especially aggregate labor supply, check out \cite*{BorellaDeNardiYang2018}.\footnote{Obviously it is easy enough to extend this further to allow for things like same-sex couples, although it requires a bit more computing power.}


\subsection{OLG Model 14: Heterogeneous Firms}
We now consider an economy with both heterogeneous households and heterogeneous firms. We can do this simply by creating the two as different permanent types of agent, one as the firm and one as the household, similarly to what was done in the previous two models. On the household side the model looks like OLG Model 6, except that we have to replace the 'asset holdings' with 'share holdings' and get rid of the progressive taxes (keeping a flat-rate earnings tax), the reason for these changes is explained below. The model contains a continuum of firms and a continuum of households, for convenience we will normalize the mass of each to 1. Note that we will have finite-lived (OLG) households, and infinitely lived firms.

We begin with the firm problem. Until now we just had a representative firm, and the reason this was possible was the assumption of constant-returns-to-scale.\footnote{Intuitively, with constant-returns-to-scale it is irrelevant if we have one big firm or two firms of half the size each of which uses half the capital and half the labor, constant-returns-to-scale means the total output remains the same, thus we can just replace all the firms with a representative. This can be proved formally.} The representative firm simply rented capital and labor from households each period. To get heterogeneous firms we will switch to firms that have decreasing-returns-to-scale production functions, and which own capital (but still rent labor).\footnote{Actually, when we had the representative firm it is irrelevant whether it was renting the capital or owning the capital, we would get the same thing in both cases \citep*{CarcelesPovedaCoenPirani2010}. We solved the representative firm as renting capital just because it was more convenient.}  Decreasing-returns-to-scale is important as now it is possible for certain allocations of capital across to be more or less efficient than others, and we will also add capital adjustment costs that mean we do not always instantly readjust to the most efficient allocations of capital across firms. The firms are a lightly modified version of \cite{GourioMiao2010}, and captures that firms might fund investment by either internal or external funds.\footnote{Internal means retained profits (or reducing dividends), external is new equity issuance; the model can be extended to distinguish debt and equity as souces of external funding but that requires an additional endogenous state variable.}

The firm problem is,
\begin{eqnarray*}
 W(k,z)&=& \max_{d,k',s,l} \left( \frac{1-\tau_d}{1-\tau_{cg}} d  -s \right)+ \frac{1}{1+r/(1-\tau_{cg})} E_{z'}[W(k',z')|z] \\
 && s + \pi = d + i + \Phi(k,i) + \tau_{corp} T \\
&& \pi = y -w l \\
&& y=z k^{\alpha_k} l^{\alpha_l} \\
&& k'= i+(1-\delta)k \\
&& \Phi(k,i)=\frac{\varphi}{2} (i/k-\delta)^2 k \\
&& T= \pi -\delta k - \phi \Phi(k,i) \\
&& d \geq 0, s \geq 0 \\
&& log(z)'=\rho_{z} log(z) + \epsilon, \quad \epsilon \sim N(0,\sigma_{\epsilon,z}^2) \\
\end{eqnarray*}
where $\pi$ is profit, $y$ is output, $k$ is physical capital, $l$ is labor, $z$ is (idiosyncratic) productivity, $d$ is dividends, $s$ is new equity issuance, $i$ is investment, $\Phi(k,i)$ is capital adjustment cost, $T$ is taxable corporate income, $\tau_{corp}$ is the corporate income tax rate. Note that the production function $y=z k^{\alpha_k} l^{\alpha_l}$ is decreasing returns to scale as long as $\alpha_k+\alpha_l<1$ (and $\alpha_k,\alpha_l>0$). The firm is getting income from new equity issuance and profits (external and internal), $s + \pi$, and uses this to cover expenses of dividends, investment, capital adjustment costs, and tax, $d + i + \Phi(k,i) + \tau_{corp} T$, labor costs being already deducted from profits. The firms differ by both their idiosyncratic productity, $z$, and their physical capital $k$.\footnote{We take advantage of a few observations to solve this problem: the first is that the decision for $l$ is actually static, and we can get a closed-form solution for it from the production function in terms of $w$, $z$ and $k$. The second is that once we choose $k'$ and one of $s$ and $d$, we can just get the other from the conditions, so we actually just have one decision variable and we will arbitrarily use $d$.}

We are now ready to discuss the reasoning for the first change we are going to make to households, namely replacing asset holdings with share holdings. Since there are heterogeneous firms we need to think about who owns them? Keeping track of the ownership between each firm and household would be prohibitively complex, and so we introduce an single 'mutual fund' which owns all firms, and households then own shares in this mutual fund. This will also have the effect of meaning that households face a single risk free return on assets (as the mutual fund will average across all the individual firm returns, and since there are a continuum of firms the return on the mutual fund will equal the mean of those individual firm returns). This is easy enough and from the household perspective mostly just looks like a renaming from asset to share, with a change at the aggregate level that we now want all shares to sum to one in equilibrium.

The household problem is,
\begin{eqnarray*}
 V(a,z,e,j)&=& \max_{c,h,a'} \frac{c^{1-\sigma}}{1-\sigma} - \psi \frac{h^{1+\eta}}{1+\eta}  + \beta s_j E_{z',e'}[V(a',j+1,z',e')|z] \\
&&\text{if working, }j<=Jr: c+P s'=(1-\tau_l) w \kappa_{j}*exp(z+e)*Lhscale* h +((1-\tau_d)D+P0)(s+AccidentBeq) - \tau_{cg}(P0-Plag)(s+AccidentBeq) \\
&&\text{if retired, }j\geq Jr: c+P s'=pension+((1-\tau_d)D+P0)(s+AccidentBeq) - \tau_{cg}(P0-Plag)(s+AccidentBeq) +\mathbb{I}_{j=J} warmglow(a') \\
&& 0\leq h \leq 1 \\
&& z'=\rho_{z} z + \epsilon, \quad \epsilon \sim N(0,\sigma_{\epsilon,z}^2) \\
&& e \sim N(0,\sigma_{e}^2)
\end{eqnarray*}
where $s$ are share holdings (in the mutual fund), $P$ is the price of mutual fund shares ($Plag$ was the price last period), the shares pay out dividends $D$, and the price of shares after paying dividends is $P0$. There is a capital gains tax $\tau_{cg}$ (on capital gain of $P0-Plag$). So the household problem looks very similar to before, except that instead of asset holdings there are share holdings, and instead of paying interest these get a dividend and a capital gain. We can compute the return to shares, denoted $r$ in terms of the dividend and share prices as,
\begin{equation*}
 1+r = \frac{P0 +(1-\tau_d)D - \tau_{cg}(P0-P)}{Plag}
\end{equation*}

Note that while I have used $z$ for shocks to both the household and firm, these are two totally different $z$ shocks (I should probably change the notation :). This is made very clear in the code where different numbers of grid points are used for each.

What is the purpose of $Lhscale$, the parameter multiplying the labour supply in the household problem? First notice that it makes no meaningful difference to the household as it is just changing the units of the labor productivity process ($\kappa_{j}*exp(z+e)*Lhscale$) which anyway had arbitrary units to begin with. The reason we need this is to ensure that the firm and household sides of the economy are the same 'size'. When we get to the general equilibrium conditions below we still have the same labor market clearance as usual. But there is no reason we should be able to get the labor supply of households to match the labor demand of firms (because both are of different arbitrary sizes; there is not just going to be a wage capable of this). So we need to think about scaling up/down the household labor supply (or firm labor demand) so that the labor market clears (at some wage). This is the purpose of $Lhscale$, we will choose it to get the labor supply of households to equal the labor demand of firms.\footnote{This is the intuition, in practice we will of course choose it joint with the other general equilibrium parameters to jointly satisfy all of the general equilibrium conditions.} But if we use $Lhscale$ to clear the labor market where does that leave the determination of the wage $w$? Notice $r$ and $w$ will determine how much firms prefer capital versus labor in production, so we will use $w$ to target a capital-labor ratio (by adding this capital-labor ratio as another general equilibrium condition).\footnote{Note that we don't need to do this same trick for the capital market clearance because of how we have the capital for firms, but just shares for households. Actually, you can look at the requirement that shares sum to one as a clever way to do exactly this for capital.}

We can now discuss why the progressive taxes have been removed from the model. Our firm is modelled as maximizing its own value.\footnote{Essentially its stock price, but in this model the stock price is purely the fundamental value of the firm.} To be able to maximize its present value it has to discount the future, and we have it do so with the discount factor $\frac{1}{1+r/(1-\tau_{cg})}$ (the firm discounts the future at the rate of return, net of capital gains tax). In principle, we want the firm to value the future at the same marginal rate of substitution of it's shareholders.\footnote{Actually, this is not so obvious. While the idea that the purpose of a firm is to maximize shareholder value is a popular concept it not clear what exactly this means if shareholders are not unanimous, nor is it clear firms actually act this way, or even that this is desirable as a reference (in Germany firms have employee representation, so clearly maximizing shareholder value is not their sole interest). Nonetheless, we here model firms as maximizing shareholder value, and have the firm discount the future at the mean marginal rate of substitution of the shareholders. We have it value the future at the mean MRS of the shareholders, but if they vote then maybe it should be the median MRS, or some other concept of the marginal shareholders MRS. All of this is even before we get to the question of whether this is actually what corporations try to do; there is a large literature on the interests of management vs shareholders, is the CEO really trying to maximize shareholder value, or do they just want to get themselves a private jet?} In the current setup, all shareholders have the same MRS, and so we don't need to calculate it.\footnote{If you take the first-order condition of the household problem you get the same intertemporal consumption Euler eqn for all unconstraint households. While the constrained households are different they are not holding any shares, and as this is the definition of the constraint we can just ignore them.} If we had progressive taxes then shareholders would have different marginal rates of substitution and so we would need to do something else, like calculating the mean of the MRS across shareholders as part of the equilibrium definition. The setup here is borrowed from \cite*{AnagnostopoulosAtesagaogluCarcelesPoveda2022}, and you can find a discussion of this issue of the MRS of shareholders in Section D on page 160-161 of \cite*{FavilukisLudvigsonVanNieuwerburgh2017} who have to deal with MRS changing due to aggregate shocks meaning their firms have a stochastic discount factor.\footnote{One other alternative is to allow firms to invest in each other, in which case they all get the market return as their own discount factor. This is described in \cite*{AnagnostopoulosAtesagaogluCarcelesPoveda2022}, and you can find the math in Section 2.3, pg 7-8, of \cite*{BarroBerkovich2018}.}$^{,\thinspace}$\footnote{We could replace the tax on labor income, $\tau_l$ with a linear tax rate on labor income and interest income. Just needs a slight change in the discount factor, see \cite*{GourioMiao2010}.}

That largely covers the setup, the rest is relatively easy, and we get the definition of the stationary general equilibrium as,
\begin{definition}
 A stationary equilibrium is prices $r$, $w$, $P0$ and government policies $\tau$ and $pension$ such that
 \begin{enumerate}
  \item Given prices ($r$, $w$, $P0$) and government policies ($\tau$, $pension$), the household value functions, $V$, and related policy function, $policy_h$, solve the household problem.
  \item Given prices ($r$, $w$, $P0$) and government policies ($\tau$, $pension$), the firm value functions, $W$, and related policy function, $policy_f$, solve the firm problem.
  \item The agents distribution of households evolves according to,
   \\  $\mu_h(a_{j+1},z_{t+1},e_{t+1},j+1)=\int \mathbb{I}_{a_{j+1}=policy_{a'}} \frac{\mu_h^{j+1}}{\mu_h^j} d\mu_h(a,z,e,j) \Gamma(z_{t+1}|z_t) Pr(\epsilon_{t+1}), \; j=1,2,...,J,\text{ and } \mu_h(:,:,:,1)=jequaloneDist$
  \item The agents distribution of firms evolves according to,
   \\ $\mu_f(k',z')=\int \mathbb{I}_{k'=policy_{k'}} d\mu_f(k,z) \Gamma(z'|z)$
  \item The aggregates are based on household policies and agent distribution,
   \\ $L_h=\int \kappa_{j} * exp(z+e)* Lhscale * policy_{h}d\mu_h$
   \\ $L_f=\int policy_l d\mu_f$
   \\ $S=\int s d\mu_h$
   \\ $Dividends=\int policy_{d} d\mu_f$
   \\ $CapitalGainsTaxRevenue=\int \tau_{cg}(P0-Plag)s d\mu_h$
   \\ $DividendTaxRevenue=(1-\tau_d) D$
   \\ $CorporateTaxRevenue=\int \tau_{corp} T d\mu_f$
  \item Labor markets clear: $L_f=L_h$
   \\ Share holdings clear: $S=1$
   \\ Dividends clear: $Dividends=D$
   \\ Pension system balance: $PensionSpending=\tau_l*w*L$
   \\ Government budget balance: $G=CapitalGainsTaxRevenue+DividendTaxRevenue+CorporateTaxRevenue$
   \\ Accidental bequest balance: $AccidentBeq=AccidentalBeqLeft/(1+n)$
 \end{enumerate}
\end{definition}
\noindent note that we do not need to include $P$ in the definition because it is captured by $r$ (given $P0$ and $D$). Because we are looking for a stationary equilibrium we also know that $Plag=P$. Note that for evolution of agent distribution of firms, the $\mu$ on both sides is the same (as this is what stationary distribution involves). Note that many of the market clearance conditions are about making sure that a variable that appears on both the households and firms side takes the same value for both, e.g., $Dividends$ is firms payment of dividends, while $D$ is household receipt of dividends.

Note that when we solve the model we are going to put the capital-labor ratio target alongside these as another general equilibrium condition (see explanation above of why we want this). Note also that we are going to choose $L_h$ as part of statationary equilibrium, I have omitted it from the definition for no particular reason.

The model that we used for heterogeneous firms is far from the only model for heterogeneous firms, in fact it is not even among the most common. If heterogeneous firms is something you are interested in you can find a few example codes (and links to materials explaining how the models work) at:
\href{https://www.vfitoolkit.com/updates-blog/2020/entry-exit-example-based-on-hopenhayn-rogerson-1993/}{vfitoolkit.com/updates-blog/2020/entry-exit-example-based-on-hopenhayn-rogerson-1993/} (there is also an example, not mentioned in link but available in \href{https://github.com/vfitoolkit/VFItoolkit-matlab-examples/tree/master/FirmEntryExitModels}{the same github repo}, based on \cite*{GourioMiao2010}, which is the model that forms the basis of the heterogeneous firms we have used here in OLG Model 14).




% \section{Search-and-Matching in OLG Models}
% \subsection{OLG Model 20: Search-and-Matching OLG}
We now solve an OLG model with search-and-matching in the labor market. From the household side this looks like an OLG model with exogenous labor, together with an infinite horizon firm problem, and a general equilibrium condition that determines the probability of finding a job (and filling a vacancy). From the perspective of VFI Toolkit, the household, firm, and general equilibrium conditions are all standard, but we need a 'custom model statistic' as one variables that goes into a general equilibrium condition is not a standard model aggregate. If you have never seen the simple Diamond-Mortensen-Pissarides search-and-matching model before, it is probably easier to look at that first before this example which combines the search-and-matching with an OLG model.

We start with the simplest search-and-matching setup. All non-employed households exert exogenous search effort. Firms create a 'vacancy' and if the match with a worker this vacancy is filled, creating a job. Workers and firms use Nash-Bargaining to determine the wage. To avoid interactions between capital and labor, it is assumed that the hiring of workers is done by intermediate firms, while production is done by a (representative) final goods firm. Thus, intermediate firms hire workers in the search-and-matching labor market, and then intermediate firms 'rent' labor to the final goods firm. The final goods firm has a standard Cobb-Douglas production function. Matches between non-employed households (who are all searching for a job) and firm vacancies are generated by a Cobb-Douglas matching function. As is standard in basic search-and-matching models, the household 'probability of finding a job' and the intermediate firm 'probability of filling a vacancy' are both proportional to the 'market tightness' which is defined as the ratio of the number of vacancies to the number of unemployed, where unemployed are the non-employed households that are searching for a job (which by assumption is here all the non-employed households, since search effort is exogenous).

Setting up search-and-matching in VFI Toolkit can be understood as involving two main aspects.

First, whether the household is employed or not is a markov state in the household problem, and the transition probability from non-employed to employed is given by the 'probability of finding a job'. Since the 'probability of finding a job' depends on the market tightness, and since the market tightness is determined in general eqm, it follows that the markov process in the household problem is determined in general eqm. To permit the markov process in the household problem to be determined in general eqm we must set it up as a function, which we do using \textit{vfoptions.ExogShockFn}.

Second, the search-and-matching adds two general eqm eqns. One of these is the markup of the wage over the marginal product of capital, which is determined by the Nash bargaining between workers and intermediate firms. This is a standard general eqm eqn so easy enough. The other general eqm eqn is the 'free entry condition' which says that the costs of (entering and) opening a vacancy must equal the benefits (otherwise more/less firms would enter to get the them equal). The challenge is that the benefits of entering depends on the value of hiring a worker, which depends on the value function of the firm and the distribution of the unemployed households (as these are the potential hires). We can set this up as a standard general eqm eqn, combined with a 'custom model statistic' for the expected value of the firm over the distribution of unemployed households. This is done using \textit{heteroagentoptions.CustomModelStats}.

To keep things simple, this model eliminates bequests and does not balance the funding for the pensions; so the model 'leaks' money due to accidental bequests that don't go to anyone, and has money 'injected' by people getting pensions for 'free'. We could add these as two more general eqm eqns, like in the other OLG models. I omit them here just so you can more easily see what search-and-matching involves in the code.

We assume that all households are either employed or non-employed, and that all non-employed households are exerting exogenous search effort. Because search-and-matching occurs randomly, and the probability of finding a job is determined in general eqm while the probability of job separation is exogenous, from the household perspective this looks like an exogenous markov process on employed/non-employed. As a result the household problem looks like a standard exogenous labor supply model, with an exogenous markov shock taking just two values, employed and non-employed. Households also have an idiosyncratic labor productivity, $z$.
\begin{eqnarray*}
 V(a,z,emp,j)&=& \max_{c,a'} \frac{c^{1-\sigma}}{1-\sigma} + \beta s_j E[V(a',z',emp',j+1)|z,emp] \\
&&\text{if working, }j<=Jr: c+a'= (1-\tau) w z emp + (1+r) a \\
&&\text{if retired, }j\geq Jr: c+a'=pension+ (1+r) a  \\
\end{eqnarray*}
Note that the interpretation of $emp$ is that $emp=1$ is employed and $emp=0$ is unemployed. We can solve this household problem in the usual manner. There is no decision variable, one endogenous state ($a$), and two exogenous markov states ($z$, $emp$).

The transition probabilities for the employment state are given by,
\begin{eqnarray*}
 Pr(emp'|emp) & emp'=0 & emp'=1 \\
 emp=0 & 1-p_{find} & p_{find} \\
 emp=1 & p_{sep} & 1-p_{sep}
\end{eqnarray*}
where $p_{sep}$ is an exogenous job seperation rate, and $p_{find}$ is a job finding probability that will be determined in general equilibrium.

There is a representative final goods firm with a Cobb-Douglas production function,
\begin{equation*}
 Y=K^{\alpha} L^{1-\alpha}
\end{equation*}
and we assume perfect competition to get the familiar conditions relating the rate of return on capital and the wage to the marginal products of capital and labor, respectively.
\begin{eqnarray*}
 h=(1-\alpha)A K^{\alpha}L^{-\alpha}
 r=\alpha A K^{\alpha-1} L^{1-\alpha} - \delta
\end{eqnarray*}
where $r$ is the return to capital, and $h$ is the rental price for labor services (which occurs in the competive market for labor services in which intermediate firms rent labor services to the final goods firm).

The labor market involves search-and-matching. Intermediate firms post vacancies, and we denote the number of vacancies as $V$. Non-employed households are all counted as unemployed (as they all exert exogenous search effort) and we denote the number of unemployed as $U$.\footnote{Techincally, $V$ is the mass of vacancies, and $U$ the mass of unemployed. Mass not number.} Matching occurs according to the Cobb-Douglas matching function,
\begin{equation*}
 M(U,V)=m U^{\gamma} V^{1-\gamma}
\end{equation*}
which leads to a job-finding probability of $p_{find}(\theta)=m \theta^{1-\gamma}$, and a job-filling probability of $q_{fill}=m \theta^{-\gamma}$, where $\theta = \frac{V}{U}$ is the market tightness.\footnote{You can easily derive this from the definitions of $p_{find}\equiv M(U,V)/U$ and $q_{fill}\equiv M(U,V)/V$. Is a direct result of the constant-returns-to-scale assumption implicit in the choice of Cobb-Douglas matching fn.}

Lastly we need to define the intermediate firm problem. Intermediate firms sit between household labor supply and final goods firm labor demand. Intermediate firms open vacancies, and then match in the labor market with households according to the job-filling probability, $q_{fill}$. Intermediate firms pay households a wage (per labor efficiency unit) of $w$, and rent this labor to the final goods firm for price $h$ (per labor efficiency unit), they make money on the difference between $h$ and $w$ (on the markup of $h$ over $w$, or to say the same thing the markdown of $w$ from $h$). Job-seperation occurs with exogenous probability $p_{sep}$.\footnote{My notation wants $q_{sep}$ to be the job-separation probility for the firm, and $p_{sep}$ to be the job-separation probability for the household, same as $p_{find}$ and $q_{fill}$ were used for job-finding and job-filling. But in equilibrium we know that $q_{sep}=p_{sep}$ (the number of jobs is the same viewed from both firm and household perspective, which is not true about job finding/filling as the vacancies and number of unemployed differ). So I just already impose the $q_{sep}=p_{sep}$ by using $p_{sep}$ directly in the intermediate firm problem.}, except that at the end of the period prior to retirement (in period $Jr-1$) the job separation rate is 1. If a worker dies, then this also results in job separation, and so the full job separation rate is $p_{sep}+sj-sj \, p_{sep}$ (assume that job seperation and death are independent shocks). So once a match occurs and a worker is hired, the value of this job to the intermediate firm is,
\begin{equation*}
 J(z,j)=(h-w) z + (1-p_{sep}-sj+sj p_{sep}) \frac{1}{1+r_{t+1}} E_z[J(z',j+1)|z]
\end{equation*}
where $J(\dot,Jr)=0$ captures the job-separation at retirement (implicitly, this is saying that $p_{sep}=1$ at age $j=Jr-1$).
notice that the per-period profit of the intermediate firm is $(h-w)z$, and then it discounts the future based on the interest rate and on the probability the job still exists next period. This leaves the issue of firms opening vacancies which lead to jobs. We assume that posting a vacancy costs $c$, and we already have that the probability of job-filling is $q_{fill}$, and that the value of the job is $J(z)$. Noting that the job occurs 'next period' and that the productivity of the worker being hired is drawn from the population of the non-employed (as these are the job seekers) we get that the free-entry condition is given by,
\begin{equation*}
 c=q_{fill} \frac{1}{1+r_{t+1}} \int_{z,j} E_z[J(z',j)|z] d\mu(emp=0,j)
\end{equation*}

The remaining issue is how the wage $w$ relates to the price of labor services $h$. We assume a Nash-inspired wage rule,
\begin{equation*}
 w=(1-\eta) \omega + \eta (h+ c \theta)
\end{equation*}
where $\eta$ is the worker's wage bargaining power.

To close the model we need to deal with the profits made by intermediate firms. We do this by distributing the profits as a lump-sum payment to households. The profit per-period made by the intermediate firms is $(h-w)z$, and so the intermediate firm sector has a profit of $\int_z (h-w)z d\mu(emp=1)$. This just requires one more simple general eqm eqn.

From the perspective of VFI Toolkit the key difference to previous models is the three additional general equilibrium conditions: (i) the differential between the wage received by households, and the wage paid by final-goods firms (which wage differential goes to the intermediate labor firms), (ii) the distribution of firm profits as a lump-sum to households, and (iii) the free-entry condition that the cost of opening a vacancy equals the expected value of the resulting job (which determines the equilibrium probability of finding a job, $p_{find}$, via market tightness).



The definition of a stationary equilibrium for this model is now,
\begin{definition}
 A stationary equilibrium is price $r$ (and $w$), and government policies $\tau$ and $pension$ such that
 \begin{enumerate}
  \item Given prices ($r$ and $w$) and government policies ($\tau$, $pension$), the household value function, $V$, and related policy function, $policy$, solve the household problem.
  \item The agents distribution evolves according to,
   \\  $\mu(a_{j+1},z_{t+1},e_{t+1},j+1)=\int \mathbb{I}_{a_{j+1}=policy_{a'}} \frac{\mu^{j+1}}{\mu^j} d\mu(a,z,e,j) \Gamma(z_{t+1}|z_t) Pr(\epsilon_{t+1}), \; j=1,2,...,J,\text{ and } \mu(:,:,:,1)=jequaloneDist$.
  \item The aggregates are based on household policies and agent distribution,
   \\ $L=\int \kappa_j exp(z+e) policy_h d\mu$
   \\ $K=\int a d\mu$
   \\ $PensionSpending=\int pension \mathbb{I}_{j\geq Jr} d\mu$
   \\ $AccidentalBeqLeft= \int policy_{a'}(1-s_j) d\hat{\mu}$
  \item Labor markets clear: $w=(1-\alpha)A K^{\alpha}L^{-\alpha}$
   \\ Capital markets clear: $r=\alpha A K^{\alpha-1} L^{1-\alpha} - \delta$
   \\ Pension system balance: $PensionSpending=\tau*w*L$
   \\ Accidental bequest balance: $AccidentBeq=AccidentalBeqLeft/(1+n)$
 \end{enumerate}
\end{definition}
\noindent note that the household problem is slightly different, and that the agent distribution (and hence all the integrals) is also now over the state space that includes the idiosyncratic exogenous shocks, namely over $(a,z,e,j)$.\footnote{When implementing we need to substitute for $Income$ with its expression in terms of the household variables (see the household problem), as otherwise the integral over the agent distribution does not make sense. I do not do so here purely because otherwise the equation gets a bit messy.} Note that the evolution of the agent distribution has been modified to also include the probabilities of the shocks.

As mentioned eariler, this model should be closed by cleaning up the accidental bequests (which currently leak out of the model) and funding the pensions (which are currently injected into the model). While the model is 'incomplete' without these they have been deliberately omitted here so that it is clearer in the codes how the search-and-matching is implemented.

Quick comment on related literature, essentially there are not many papers doing this. De la Croix, Pierrard \& Sneessens (2013) - Aging and pensions in general equilibrium: Labor market imperfections matter, does an OLG with search-matching, but they use one big representative household so not really relevant. Lugauer (2012) - Demographic Change and the Great Moderation in an Overlapping Generations Model with Matching Frictions, is one of the first examples. Cheron, Hairault \& Langot (2013) - Life-Cycle Equilibrium Unemployment, is the other. Very few other papers seem to do this, Graham \& Ozbilgin (2021) - Age, industry, and unemployment risk during a pandemic lockdown, does transition path in an OLG with search-and-matching, plus they have endogenous separation. What we are doing here is random search, Menzio, Telyukova \& Visschers (2016) - Directed search over the life cycle, does something similar but with directed search.

\subsection{OLG Model 21: Endogenous Search Effort}
We make one change to OLG Model 20, namely endogenizing search effort. This leads to one major conceptual change, namely the model now distinguishes employed, unemployed and inactive (previously, there was just employed and non-employed); with unemployed being those non-employed and actively searching (search effort greater than zero) and inactive being those non-employed but not searching (search effort zero). To keep the model tractable, it is assumed that the probability of finding a job is linear in the search effort (doubling search effort doubles the probability of finding a job).

In terms of what this means for setting up and solving the model, there are some minor changes to general equilibrium conditions and the definitions of parameters, but the major change is to the household problem. The household problem now includes the search effort as a decision and that this decision influences the transitions to employment which in turn means that the employed/non-employed variable is now a semi-exogenous state (not an exogenous markov state as before).

The household problem is now,
\begin{eqnarray*}
 V(a,semiz,emp,j)&=& \max_{d,c,a'} \frac{c^{1-\sigma}}{1-\sigma}-\psi(d) + \beta s_j E[V(a',semiz',emp',j+1)|z,emp] \\
&&\text{if working, }j<=Jr: c+a'= (1-\tau) w z emp + (1+r) a \\
&&\text{if retired, }j\geq Jr: c+a'=pension+ (1+r) a  \\
&& semiz'(d,semiz)
\end{eqnarray*}
where $d$ is assumed to take three possible values (0,1,2).

The transition probabilities for the employment semi-exogenous state are given by,
\begin{eqnarray*}
 Pr(emp'|emp) & emp'=0 & emp'=1 \\
 emp=0 & 1-d\,p_{find} & d\,p_{find} \\
 emp=1 & p_{sep} & 1-p_{sep}
\end{eqnarray*}
where $p_{sep}$ is an exogenous job seperation rate, and $d\, p_{find}$ is a job finding probability. $p_{find}$ is the probability of finding a job per unit of search effort, and will be determined in general equilibrium.

The final goods firm problem is unchanged. The only change the the matching function in the search-and-matching labor market is that $U$ is now defined as $\int d (emp=0) d\mu$ (previously it was $\int (emp=0) d\mu$). The problem of the intermediate firm is unchanged, except that the free-entry condition is now,
\begin{equation*}
 c=q_{fill} \frac{1}{1+r_{t+1}} \int_{z,j} E_z[J(z',j)|z] d\, d\mu(emp=0,j)
\end{equation*}
where we the integral with respect to the non-employed is now weighted by their search effort. The determination of wages is unchanged.






\newpage

\newpage

\bibliographystyle{plainnat}
\bibliography{/home/kmarshallbanana/Dropbox/RDKbiblio.bib}

\newpage
\appendix

\section{Stationary General Equilibrium: Explanation of Solution Method}
What is happening when we solve for general equilibrium? The basic idea is that we are looking for the parameters (specifically those parameters that are being determined in general equilibrium) that make the general equilibrium conditions all equal to zero. Recall that when we write the general equilibrium conditions in the code we write them as expressions that will evaluate to zero when the condition holds.

We can think of this as choosing the general equilibrium parameters, $\theta$, to solve following optimization problem,
\begin{eqnarray*}
 \min_{\theta} \sum_g GEcond_g(\theta)^2
\end{eqnarray*}
where the $\sum_g GEcond_g(\theta)^2$ is the sum of squares of the general equilibrium conditions (note that by squaring a variable we make it so that the minimum possible value is zero; think about what a graph of $x^2$ looks like).\footnote{You could alternatively take the absolute value of each general equilibrium condition (actually, just set $heteroagentoptions.multiGEcriterion=0$ and it will do so). Also note that we are essentially giving every general equilibrium condition a weight of one and then taking the weighted sum. You can modify these weights using $heteroagentoptions.multiGEweights$.} Minimizing this sum of squares is going to result in finding $\theta$ such that all the general equilibrium conditions equal zero (if possible).

When using the code $\theta$ is the parameters we name in \textit{GEPriceParamNames}, and the general equilibrium conditions are those we specify in $GeneralEqmEqns$ which should always be written to evaluate to zero in general equilibrium.

So, finding the general equilibrium parameters such that the all the general equilibrium conditions equal zero (which will be the minimum of the sum of squares of the general equilibrium conditions) is the problem that the general equilibrium commands are solving, but how do they do this?

The command to find the general equilibrium is doing essentially the following steps,
\begin{enumerate}
 \item Guess general equilibrium parameters $\theta$.
 \item Based on $\theta$, solve the value function problem to get the policy function.
 \item Based on policy function, solve for the agent stationary distribution.
 \item Based on $\theta$, policy and agent stationary distribution, evaluate the FnsToEvaluate.
 \item Based on $\theta$ and evaluated FnsToEvaluate, evaluate the general equilibrium conditions.
 \item If the general equilibrium conditions are all zero, finish (and report $\theta$). Otherwise, update $\theta$ to a new guess and go back to step 1.
\end{enumerate}
Obviously the word 'update' is hiding a lot of detail about how to update $\theta$ in an intelligent fashion.\footnote{There is a vast computational literature on how to solve optimization problems, such as minimization. Appendix B here gives a very brief descriptions of some alternative options to the default in the codes.} But this is enough for us to understand the main steps.


\section{Stationary General Equilibrium: How to Solve Faster, Or solve more robustly}
By default VFI Toolkit uses Matlab's \textit{fminsearch()} command to solve the minimization problem of choosing the general equilibrium parameters to minimize (the sum of squares of) the general equilibrium conditions. This is easy to use, and provides a decent mix of speed and robustness for a default setting. In this section we explain three alternatives to the default of \textit{fminsearch()}, the first is to use Matlab's non-linear least squares optimizer \textit{lsqnonlin()} which is easy and fast but requires the Matlab Optimization Toolbox and as it is derivatives-based can be less accurate, the second is to use a shooting algorithm which is faster but requires more setting up and is less robust, and the third is to use the CMA-ES algorithm which is slower but much more robust and is trivial to set up. Code for OLGModelB implements these for OLG Model 6. Another option is to use to optimization algorithms one after another, which the toolkit can do for you.

Relatedly, you can set the accuracy that is used for the general equilibrium equations using \textit{heteroagentoptions.toleranceGEcondns=10ˆ(-4)}, and of the general equilibrium parameters using \textit{heteroagentoptions.toleranceGEprices=10ˆ(-4)} (these will be passed internally to the optimization routines to use as their 'stopping tolerance').

Of course one other easy trick for making things faster is to first make the grids small, solve the general equilibrium which will be fast because of the small grids, and then use this as an initial guess for the general equilibrium parameters while solving for a larger grid (this is known as a multi-grid method).

We will start out with \textit{lsqnonlin()}. Least-squares minimization problems have specific structure that can be exploited during solution to speed things up, and \textit{lsqnonlin()} does this in a way that depends on derivatives. This makes it fast and powerful, but also means that sometimes it only solves the general eqm eqns to about three decimal places. Using \textit{lsqnonlin()} is as simple as setting \textit{heteroagentoptions.fminalgo=8}.

Now we see how to use CMA-ES (covariance-matrix adaptation - evolutionary strategy). We do not attempt to explain how CMA-ES works here, you can search online where there is pleny of material. Actually implementing in codes to use CMA-ES instead of \textit{fminsearch()} is as simple as setting \textit{heteroagentoptions.fminalgo=4}. Done. (heteroagentoptions.fminalgo=1 is the default of fminsearch). This will make codes slower but has a much better ability to find the general equilibrium in complicated problems and where the initial guess is bad (loosely, it is going to spend more time moving slowly around the parameter space to find the optimum, and is less likely to get stuck at local minima).

Setting up the shooting algorithm takes a little more work, but reduces run times notably (half or less). The idea is roughly as follows. We often can think how to update a given price based on a specific general equilibrium conditions. For example, in a standard model with a representative firm with Cobb-Douglas production function and perfect competition in the capital market, the interest rate $r$ is going to equal the marginal product of capital (net of depreciation), that is $r=\alpha K^{\alpha-1} L^{1-\alpha}-\delta$. So we have the general equilibrium conditions $r-\alpha K^{\alpha-1} L^{1-\alpha}-\delta$ (general equilibrium being when this equals zero). Notice that if this is positive, then $r$ is to big so we could reduce $r$, say by subtracting the value of this condition, and if this general equilibrium condition is negative then $r$ is to small and we could increase $r$ again by subtracting the value of this condition (which is negative). So maybe we could write something like $r_{new}=r_{old}-\phi (r_{old}-\alpha K^{\alpha-1} L^{1-\alpha}-\delta)$, that is the new price is the old price minus coefficient times general equilibrium condition. This is a shooting algorithm. To set it up we first tell the code we want to use the shooting algorithm, and then tell it which general equilibrium condition to use to update which price, whether to add or subtract, and what value to give the coefficient. If you use bigger coefficients the code will be faster but more fragile (might fail), if you use smaller coefficients the code will be slower but more stable. First, to say we want to use the shooting algorithm
\begin{equation*}
 heteroagentoptions.fminalgo=5;
\end{equation*}
then we define how to update the general equilibrium prices based on the general equilibrium conditions,
\begin{table}[h!]
\footnotesize
heteroagentoptions.fminalgo5.howtoupdate=\{...  \% a row is: GEcondn, price, add, factor \\
\phantom{aaa} 'capitalmarket','r',0,0.1;...  \% capitalmarket GE condition will be positive if r is too big, so subtract \\
\phantom{aaa} 'pensions','pension',0,0.1;... \% pensions GE condition will be positive if pension is too big, so subtract \\
\phantom{aaa} 'bequests','AccidentBeq',1,0.1;... \% bequests GE condition will be positive if AccidentBeq is too small, so add \\
\phantom{aaa} 'Gtarget','G',0,0.1;... \% Gtarget GE condition will be positive if G is too big, so subtract \\
\phantom{aaa} 'govbudget','eta1',1,0.1;... \% govbudget GE condition will be negative if eta1 is too big, so add \\
    \};
\end{table}
To make sense of these you need to look at the general equilibrium conditions themselves (this is OLG Model 6). Note: the update is essentially new\_price=price+factor*add*GEcondn\_value-factor*(1-add)*GEcondn\_value. Notice that this adds factor*GEcondn\_value when add=1 and subtracts it what add=0. A small 'factor' will make the convergence to solution take longer, but too large a value will make it  unstable (fail to converge). Technically this is the damping factor in a shooting algorithm. There is no reason why I use 0.1, nor why this is the same for all of them. You can try changing a few to 0.2 or even 0.5 to see how the code runs faster (or if too high might even fail to solve).

Lastly, we can combine these optimization methods. For example we know that \textit{lsqnonlin()} (\textit{heteroagentoptions.fminalgo=8}) is fast, but that because it is based on derivatives can struggle to reach high levels of accuracy, so we might want to use it but then clean up the solution with \textit{fminsearch()} (\textit{heteroagentoptions.fminalgo=1}). Just setting \textit{heteroagentoptions.fminalgo=[8,1]}, the toolkit will automatically first use \textit{lsqnonlin()}, and then starting from the solution this gives it switches to \textit{fminsearch()}. While using this, you might also want to set that the accuracy you require of the first method is less than the final accuracy, so you could, e.g., set \textit{heteroagentoptions.toleranceGEcondns=[10ˆ(-3),10ˆ(-4)]} which means that it will use the first of these with \textit{lsqnonlin()}, and the second with \textit{fminsearch()}.

\section{Stationary General Equilibrium: Constrain General Equilibrium Parameters}
Say you are finding the wage, $w$, in general equilibrium. You know this has to be positive, and if you try to solve the model with a negative wage then it gives errors. You can set constraints on the general equilibrium parameters to deal with these situations. There are three possible constraints, all of which are applied by specifying the name of the parameter you want to constrain. In the following, we will use $w$ as the name of the parameter (in GEPriceParamNames) that we wish to constain.

You can constrain a parameter to be positive using
\\ \noindent \textit{heteroagentoptions.constrainpositive=\{'w'\}}
You can constrain a parameter to be between 0 and 1 (e.g., because it is a probability) using
\\ \noindent \textit{heteroagentoptions.constrain0to1=\{'w'\}}
You can constrain a parameter to be between A and B using
\\ \noindent \textit{heteroagentoptions.constrainAtoB=\{'w'\}}
and then you specify the values of A and B, e.g. that $w$ is between 3 and 5, with
\\ \noindent \textit{heteroagentoptions.constrainAtoBlimits.w=[3,5]}
[Use name of parmeter as field, '.w', and then use a vector of two values for the range to apply to that parameter]



\section{General Equilibrium beyond Stationary General Equilibrium} \label{app:eqmwithgrowth}
Not yet written. Will explain a wider definition of general equilibrium, and then explain the definition of stationary general equilibrium. Is not something any of the models here actually do anyway.







\end{document}
